
\PassOptionsToPackage{draft}{hyperref}
\documentclass[12pt,a4paper]{article}

% Compilar con LuaLaTeX + Biber.
\usepackage{viu-mrob-thesis}
\usepackage{csquotes}
\usepackage[
  backend=biber,
  style=apa,
  natbib=true,
  hyperref=false,
  sorting=nyt,
  sortcites=true
]{biblatex}
\addbibresource{references.bib}

\usepackage{algorithm}
\usepackage{algpseudocode}
\usepackage[export]{adjustbox}
\usepackage{subcaption}
\usepackage{array}
\usepackage{tabularx}
\usepackage{needspace}

\floatname{algorithm}{Algoritmo}
\renewcommand{\algorithmicrequire}{\textbf{\textcolor{VIUOrange}{Entrada:}}}
\renewcommand{\algorithmicensure}{\textbf{\textcolor{VIUOrange}{Salida:}}}
\algrenewcommand\algorithmicif{\textcolor{VIUOrange}{\textbf{si}}}
\algrenewcommand\algorithmicthen{\textcolor{VIUOrange}{\textbf{entonces}}}
\algrenewcommand\algorithmicelse{\textcolor{VIUOrange}{\textbf{si no}}}
\algrenewcommand\algorithmicfor{\textcolor{VIUOrange}{\textbf{para}}}
\algrenewcommand\algorithmicforall{\textcolor{VIUOrange}{\textbf{para cada}}}
\algrenewcommand\algorithmicreturn{\textcolor{VIUOrange}{\textbf{devolver}}}
\algrenewcommand\algorithmicfunction{\textcolor{VIUOrange}{\textbf{función}}}
\algrenewcommand\alglinenumber[1]{\scriptsize\textcolor{VIUGray}{#1:}}
\algrenewcommand\algorithmiccomment[1]{\hfill\textcolor{VIUGray}{\itshape\% #1}}
\algtext*{EndIf}
\algtext*{EndFor}
\algtext*{EndFunction}

\makeatletter
\renewcommand{\ALG@beginalgorithmic}{\fontsize{9}{10.5}\selectfont}
\makeatother

% Generated by scripts/build_sp1_levels_pdf.py; do not edit.
\newcommand{\NOneRows}{12\,630}
\newcommand{\NOneMaxN}{2\,048}
\newcommand{\NOneExponent}{2{,}40}
\newcommand{\NOneRSquared}{0{,}997}
\newcommand{\NOneGreedyRatio}{1{,}070}
\newcommand{\NOneWorlds}{6\,630}
\newcommand{\NOneQualityRows}{4\,500}
\newcommand{\NOneSavingPct}{6{,}53}
\newcommand{\NOneSavingLowPct}{6{,}34}
\newcommand{\NOneSavingHighPct}{6{,}71}
\newcommand{\NOneScenarioGates}{3}
\newcommand{\NOneCellsPerScenario}{9}
\newcommand{\NOneUniformCellsAbove}{9}
\newcommand{\NOneClusteredCellsAbove}{9}
\newcommand{\NOneSeparatedCellsAbove}{0}
\newcommand{\NOneRingCellsAbove}{0}
\newcommand{\NOneCorridorCellsAbove}{8}
\newcommand{\NOneCorridorSmallExtremePct}{4{,}75}
\newcommand{\NOneQualityPHolmBound}{4{,}70 \times 10^{-108}}
\newcommand{\NOneUniformSavingPct}{11{,}38}
\newcommand{\NOneClusteredSavingPct}{9{,}92}
\newcommand{\NOneSeparatedSavingPct}{1{,}18}
\newcommand{\NOneRingSavingPct}{1{,}66}
\newcommand{\NOneCorridorSavingPct}{8{,}27}
\newcommand{\NOneExponentLow}{2{,}37}
\newcommand{\NOneExponentHigh}{2{,}44}
\newcommand{\NOneScalingPninetyfiveMs}{745{,}9}
\newcommand{\NOneScalingCompleted}{630}
\newcommand{\NOneScalingFailed}{0}
\newcommand{\NOneProcessor}{Intel Core Ultra 9 185H}
\newcommand{\NOneLogicalCPU}{22}
\newcommand{\NOnePythonVersion}{3.13.14}
\newcommand{\NOneSciPyVersion}{1.16.3}
\newcommand{\NOneFailureRows}{6\,000}
\newcommand{\NOneFailureWorlds}{1\,200}
\newcommand{\NOneFailureAgreementPct}{100}
\newcommand{\NOneHeteroRows}{1\,500}
\newcommand{\NOneHeteroWorlds}{300}
\newcommand{\NOneLowFalsePct}{38}
\newcommand{\NOneModerateFalsePct}{72}
\newcommand{\NOneHighFalsePct}{92{,}3}
\newcommand{\NOneExtremeFalsePct}{97{,}7}
\newcommand{\NOneMilpCertifiedPct}{99{,}60}
\newcommand{\NOneMilpAudits}{1\,500}
\newcommand{\NOneMilpFeasibleIncumbents}{1\,500}
\newcommand{\NOneMilpCertifiedCount}{1\,494}
\newcommand{\NOneMilpUncertifiedCount}{6}
\newcommand{\NOneFalseFeasibleCount}{900}
\newcommand{\NOneMilpFeasibleAmongFalse}{900}
\newcommand{\NOneMilpCertifiedAmongFalse}{897}
\newcommand{\NOneMilpUncertifiedFalseCount}{3}
\newcommand{\NOneMilpGapMinPct}{0{,}83}
\newcommand{\NOneMilpGapMaxPct}{2{,}14}
\newcommand{\NOneMilpTimeLimitSec}{5}
\newcommand{\NOneHeteroPHolmBound}{4{,}81 \times 10^{-35}}
\newcommand{\NOneHeteroContrastsSupported}{4}
\newcommand{\NOneShareUniform}{90{,}4}
\newcommand{\NOneShareClustered}{96{,}7}
\newcommand{\NOneShareCorridor}{85{,}1}
\newcommand{\NOneShareSeparated}{0{,}1}
\newcommand{\NOneShareRing}{3{,}9}
\newcommand{\NOneOrderOverallPct}{3{,}76}
\newcommand{\NOneOrderUniformPct}{6{,}84}
\newcommand{\NOneOrderCorridorPct}{5{,}31}
\newcommand{\NOneOrderCorridorLowPct}{5{,}14}
\newcommand{\NOneFailureCostSharePct}{99{,}7}
\newcommand{\NOneFailureCostRows}{3\,000}
\newcommand{\NOneMaxTestedFailureFraction}{0{,}30}
\newcommand{\NOneFailureCostMaxPct}{41{,}5}
\newcommand{\NOneDeficitLowPct}{4{,}7}
\newcommand{\NOneDeficitModeratePct}{14{,}3}
\newcommand{\NOneDeficitHighPct}{33{,}2}
\newcommand{\NOneDeficitExtremePct}{53{,}9}
\newcommand{\NOneDeficitExtremeLowPct}{50{,}5}
\newcommand{\NOneDeficitExtremeHighPct}{57{,}6}
\newcommand{\NOneTrendSlope}{8{,}58}
\newcommand{\NOneTrendLow}{7{,}46}
\newcommand{\NOneTrendHigh}{9{,}69}
\newcommand{\NOneTrendOdds}{2{,}36}
\newcommand{\NOneTrendClusters}{300}
\newcommand{\NOneTiebreakAgreementPct}{100}
\newcommand{\NOneCrossLoadTies}{0}
\newcommand{\NOneCertificateWorlds}{300}
\newcommand{\NOneCertificateFalse}{0}
\newcommand{\NOneUncertifiedWorlds}{1\,200}
\newcommand{\NOneUncertifiedFalse}{900}
\newcommand{\NTwoRows}{4\,480}
\newcommand{\NTwoOracleRows}{750}
\newcommand{\NTwoOracleAgreementPct}{100}
\newcommand{\NTwoOracleDisagreements}{0}
\newcommand{\NTwoOracleComparable}{748}
\newcommand{\NTwoOracleInfeasible}{2}
\newcommand{\NTwoOracleMaxError}{4{,}15 \times 10^{-12}}
\newcommand{\NTwoOracleMaxStates}{4\,194\,304}
\newcommand{\NTwoHomogeneousRows}{250}
\newcommand{\NTwoHomogeneousAgreementPct}{100}
\newcommand{\NTwoHomogeneousMaxError}{4{,}55 \times 10^{-13}}
\newcommand{\NTwoHomogeneousCardinality}{250}
\newcommand{\NTwoAtomicityRows}{1\,500}
\newcommand{\NTwoAtomicityComparable}{1\,455}
\newcommand{\NTwoGapPct}{18}
\newcommand{\NTwoGapLowPct}{17{,}6}
\newcommand{\NTwoGapHighPct}{18{,}5}
\newcommand{\NTwoGapPninetyfivePct}{35{,}4}
\newcommand{\NTwoGapMaxPct}{74}
\newcommand{\NTwoGapMinPct}{1{,}7}
\newcommand{\NTwoFractionalPct}{100}
\newcommand{\NTwoFractionalLowPct}{99{,}7}
\newcommand{\NTwoNegativeGaps}{0}
\newcommand{\NTwoStrictGapRows}{1\,455}
\newcommand{\NTwoLpOnlyCount}{45}
\newcommand{\NTwoLpOptimalRows}{1\,500}
\newcommand{\NTwoOtherStatusRows}{0}
\newcommand{\NTwoFractionalAllPct}{100}
\newcommand{\NTwoSpearman}{0{,}082}
\newcommand{\NTwoGapHomogeneousPct}{23{,}1}
\newcommand{\NTwoGapMidPct}{13{,}9}
\newcommand{\NTwoGapExtremePct}{22{,}8}
\newcommand{\NTwoPhaseRows}{1\,800}
\newcommand{\NTwoPhaseWorlds}{360}
\newcommand{\NTwoPhaseCensored}{29}
\newcommand{\NTwoSlackRows}{1\,500}
\newcommand{\NTwoSlackInfeasible}{163}
\newcommand{\NTwoSlackInfeasiblePct}{10{,}9}
\newcommand{\NTwoTrendCv}{3{,}30}
\newcommand{\NTwoTrendCvLow}{2{,}91}
\newcommand{\NTwoTrendCvHigh}{3{,}69}
\newcommand{\NTwoTrendPressure}{-35{,}4}
\newcommand{\NTwoTrendPressureLow}{-40{,}2}
\newcommand{\NTwoTrendPressureHigh}{-30{,}6}
\newcommand{\NTwoPhaseHomogeneousPct}{41{,}9}
\newcommand{\NTwoPhaseExtremePct}{81{,}1}
\newcommand{\NTwoBandHomogeneousPct}{17{,}2}
\newcommand{\NTwoBandMiddlePct}{100}
\newcommand{\NTwoBandExtremePct}{95{,}6}
\newcommand{\NTwoBandPairs}{180}
\newcommand{\NTwoBandMiddleCv}{0{,}35}
\newcommand{\NTwoBandMiddleOnly}{8}
\newcommand{\NTwoBandExtremeOnly}{0}
\newcommand{\NTwoBandMcNemarP}{0{,}0078}
\newcommand{\NTwoPlateauLow}{0{,}35}
\newcommand{\NTwoPlateauHigh}{0{,}65}
\newcommand{\NTwoCertRows}{180}
\newcommand{\NTwoCertMinN}{20}
\newcommand{\NTwoCertMaxN}{80}
\newcommand{\NTwoIncumbentPct}{100}
\newcommand{\NTwoCertNTwenty}{100}
\newcommand{\NTwoCertNForty}{93{,}3}
\newcommand{\NTwoCertNFifty}{66{,}7}
\newcommand{\NTwoCertNSixty}{33{,}3}
\newcommand{\NTwoCertNEighty}{26{,}7}
\newcommand{\NTwoCensoredRows}{85}
\newcommand{\NTwoGapCertMinPct}{0{,}4}
\newcommand{\NTwoGapCertMedPct}{2{,}8}
\newcommand{\NTwoGapCertPninetyfivePct}{8{,}8}
\newcommand{\NTwoGapCertMaxPct}{11{,}7}
\newcommand{\NThreeWorlds}{1\,200}
\newcommand{\NThreeRows}{3\,600}
\newcommand{\NThreeOracleFeasible}{1\,159}
\newcommand{\NThreeOracleInfeasible}{41}
\newcommand{\NThreeOracleCertified}{1\,159}
\newcommand{\NThreeFalseFeasible}{0}
\newcommand{\NThreeCommonSupport}{878}
\newcommand{\NThreeCochranP}{6{,}7 \times 10^{-122}}
\newcommand{\NThreeFriedmanP}{4{,}5 \times 10^{-254}}
\newcommand{\NThreePartitionComponents}{2}
\newcommand{\NThreeControlRateMaxPct}{28}
\newcommand{\NThreeConnectedRateMinPct}{54}
\newcommand{\NThreeCyclesObserved}{1\,200}
\newcommand{\NThreeInconsistent}{0}
\newcommand{\NThreeCochranQ}{558}
\newcommand{\NThreeFriedmanQ}{1166{,}7}
\newcommand{\NThreeKendallW}{0{,}664}
\newcommand{\NThreeMcNemarP}{4{,}1 \times 10^{-84}}
\newcommand{\NThreeMcNemarPairP}{0{,}50}
\newcommand{\NThreeRatioBytesGrape}{2{,}30}
\newcommand{\NThreeRatioBytesPair}{2{,}84}
\newcommand{\NThreeRatioMsgGrape}{8{,}6}
\newcommand{\NThreeRatioMsgPair}{10{,}6}
\newcommand{\NThreeGrapeCertAgree}{300}
\newcommand{\NThreeGrapeCertWorlds}{300}
\newcommand{\NThreeGrapeCostIdentical}{300}
\newcommand{\NThreePairCostIdentical}{38}
\newcommand{\NThreePairMaxSpread}{221{,}1}
\newcommand{\NThreeGapCompletePair}{10{,}5}
\newcommand{\NThreeGapCompleteGrape}{31}
\newcommand{\NThreeGapDensePair}{11{,}9}
\newcommand{\NThreeGapDenseGrape}{31}
\newcommand{\NThreeGapMediumPair}{14{,}8}
\newcommand{\NThreeGapMediumGrape}{31}
\newcommand{\NThreeGapThresholdPair}{18{,}9}
\newcommand{\NThreeGapThresholdGrape}{31}
\newcommand{\NThreeDiffGrapeCbbaPct}{24{,}07}
\newcommand{\NThreeDiffGrapeCbbaLowPct}{21{,}57}
\newcommand{\NThreeDiffGrapeCbbaHighPct}{26{,}49}
\newcommand{\NThreeDiffPairCbbaPct}{24{,}25}
\newcommand{\NThreeDiffPairCbbaLowPct}{21{,}74}
\newcommand{\NThreeDiffPairCbbaHighPct}{26{,}75}
\newcommand{\NThreeDiffPairGrapePct}{0{,}17}
\newcommand{\NThreeDiffPairGrapeLowPct}{0{,}00}
\newcommand{\NThreeDiffPairGrapeHighPct}{0{,}43}
\newcommand{\NThreeFeasPctCbba}{75{,}8}
\newcommand{\NThreeGapPctCbba}{40{,}3}
\newcommand{\NThreeCommonGapPctCbba}{40{,}3}
\newcommand{\NThreeBytesCbba}{9\,089}
\newcommand{\NThreeRoundsCbba}{25}
\newcommand{\NThreeCompletePctCbba}{54}
\newcommand{\NThreeMediumPctCbba}{61{,}3}
\newcommand{\NThreeThresholdPctCbba}{66{,}7}
\newcommand{\NThreePartitionedPctCbba}{13{,}3}
\newcommand{\NThreeFeasPctGrape}{99{,}8}
\newcommand{\NThreeGapPctGrape}{22{,}1}
\newcommand{\NThreeCommonGapPctGrape}{20{,}9}
\newcommand{\NThreeBytesGrape}{20\,909}
\newcommand{\NThreeRoundsGrape}{225}
\newcommand{\NThreeCompletePctGrape}{99{,}3}
\newcommand{\NThreeMediumPctGrape}{99{,}3}
\newcommand{\NThreeThresholdPctGrape}{99{,}3}
\newcommand{\NThreePartitionedPctGrape}{22}
\newcommand{\NThreeFeasPctPair}{100}
\newcommand{\NThreeGapPctPair}{10{,}4}
\newcommand{\NThreeCommonGapPctPair}{9{,}4}
\newcommand{\NThreeBytesPair}{25\,791}
\newcommand{\NThreeRoundsPair}{289}
\newcommand{\NThreeCompletePctPair}{100}
\newcommand{\NThreeMediumPctPair}{100}
\newcommand{\NThreeThresholdPctPair}{99{,}7}
\newcommand{\NThreePartitionedPctPair}{28}
\newcommand{\NFourWorlds}{1\,200}
\newcommand{\NFourRows}{12\,000}
\newcommand{\NFourCThreeGapDiffPct}{-49{,}2}
\newcommand{\NFourCThreeGapDiffLowPct}{-53{,}1}
\newcommand{\NFourCThreeGapDiffHighPct}{-45{,}9}
\newcommand{\NFourCThreeGapP}{9{,}9 \times 10^{-187}}
\newcommand{\NFourCThreeRiskPct}{2{,}33}
\newcommand{\NFourCThreeRiskLowPct}{1{,}47}
\newcommand{\NFourCThreeRiskHighPct}{3{,}19}
\newcommand{\NFourSmithGapDiffPct}{9{,}7}
\newcommand{\NFourLllGapDiffPct}{6{,}8}
\newcommand{\NFourDGapDiffPct}{0{,}00}
\newcommand{\NFourDRiskPct}{0{,}00}
\newcommand{\NFourDByteDiff}{1\,282}
\newcommand{\NFourDByteDiffLow}{1\,258}
\newcommand{\NFourDByteDiffHigh}{1\,318}
\newcommand{\NFourDRuntimeDiff}{-26{,}1}
\newcommand{\NFourDpopWorlds}{400}
\newcommand{\NFourDpopMatches}{400}
\newcommand{\NFourDpopMaxError}{1{,}73 \times 10^{-12}}
\newcommand{\NFourDpopMaxEntries}{3\,279}
\newcommand{\NFourDpopMaxBytes}{26\,288}
\newcommand{\NFourScaleSlopeD}{3{,}79}
\newcommand{\NFourScaleSlopeDLow}{3{,}46}
\newcommand{\NFourScaleSlopeDHigh}{4{,}11}
\newcommand{\NFourScaleRSquaredD}{0{,}999}
\newcommand{\NFourScaleSlopeU}{1{,}89}
\newcommand{\NFourScaleSlopeULow}{1{,}59}
\newcommand{\NFourScaleSlopeUHigh}{2{,}19}
\newcommand{\NFourScaleRSquaredU}{0{,}997}
\newcommand{\NFourScaleRuntimeDMax}{20614{,}8}
\newcommand{\NFourScaleRuntimeUMax}{99{,}7}
\newcommand{\NFourScaleBytesDMax}{52\,097}
\newcommand{\NFourScaleBytesUMax}{2\,527}
\newcommand{\NFourPopLogitHundredPct}{100}
\newcommand{\NFourPopBnnHundredPct}{10}
\newcommand{\NFourPopRepHundredPct}{0}
\newcommand{\NFourPopSmithHundredPct}{0}
\newcommand{\NFourFeasPctCbba}{79}
\newcommand{\NFourGapPctCbba}{42{,}2}
\newcommand{\NFourBytesCbba}{9\,188}
\newcommand{\NFourRoundsCbba}{25}
\newcommand{\NFourRuntimeCbba}{26{,}5}
\newcommand{\NFourFeasPctGrape}{99{,}7}
\newcommand{\NFourGapPctGrape}{21{,}6}
\newcommand{\NFourBytesGrape}{20\,608}
\newcommand{\NFourRoundsGrape}{225}
\newcommand{\NFourRuntimeGrape}{121{,}5}
\newcommand{\NFourFeasPctPairGrape}{100}
\newcommand{\NFourGapPctPairGrape}{9{,}7}
\newcommand{\NFourBytesPairGrape}{25\,482}
\newcommand{\NFourRoundsPairGrape}{289}
\newcommand{\NFourRuntimePairGrape}{293{,}1}
\newcommand{\NFourFeasPctQpgU}{97{,}7}
\newcommand{\NFourGapPctQpgU}{55{,}7}
\newcommand{\NFourBytesQpgU}{662}
\newcommand{\NFourRoundsQpgU}{39}
\newcommand{\NFourRuntimeQpgU}{7{,}9}
\newcommand{\NFourFeasPctSmith}{96{,}7}
\newcommand{\NFourGapPctSmith}{68{,}4}
\newcommand{\NFourBytesSmith}{658}
\newcommand{\NFourRoundsSmith}{48}
\newcommand{\NFourRuntimeSmith}{8}
\newcommand{\NFourFeasPctLll}{98{,}2}
\newcommand{\NFourGapPctLll}{65{,}9}
\newcommand{\NFourBytesLll}{6\,084}
\newcommand{\NFourRoundsLll}{155}
\newcommand{\NFourRuntimeLll}{30{,}6}
\newcommand{\NFourFeasPctQpgP}{99{,}9}
\newcommand{\NFourGapPctQpgP}{14{,}8}
\newcommand{\NFourBytesQpgP}{3\,324}
\newcommand{\NFourRoundsQpgP}{281}
\newcommand{\NFourRuntimeQpgP}{142{,}6}
\newcommand{\NFourFeasPctCThree}{100}
\newcommand{\NFourGapPctCThree}{3{,}2}
\newcommand{\NFourBytesCThree}{14\,164}
\newcommand{\NFourRoundsCThree}{845}
\newcommand{\NFourRuntimeCThree}{250{,}5}
\newcommand{\NFourFeasPctQpgCf}{100}
\newcommand{\NFourGapPctQpgCf}{11{,}6}
\newcommand{\NFourBytesQpgCf}{4\,080}
\newcommand{\NFourRoundsQpgCf}{16}
\newcommand{\NFourRuntimeQpgCf}{185{,}2}
\newcommand{\NFourFeasPctQpgD}{100}
\newcommand{\NFourGapPctQpgD}{11{,}6}
\newcommand{\NFourBytesQpgD}{5\,355}
\newcommand{\NFourRoundsQpgD}{16}
\newcommand{\NFourRuntimeQpgD}{155{,}9}

\InputIfFileExists{generated/n4_v3_results_macros.tex}{}{}
\InputIfFileExists{generated/n4_v4_results_macros.tex}{}{}
\InputIfFileExists{generated/sp1-final/review_counts.tex}{}{}
\InputIfFileExists{generated/sp1-final/canonical_metrics.tex}{}{}
\InputIfFileExists{generated/sp1-final/traceable_counts.tex}{}{}
\InputIfFileExists{generated/sp1-final/new_campaigns.tex}{}{}

\hypersetup{
  pdftitle={Resultados SP1 compactados},
  pdfauthor={Jorge Luis Mayorga Taborda},
  pdfsubject={SP1: formación de coaliciones de AMR}
}

% --------------------------------------------------------------------------
% Ajustes editoriales: 12 pt en cuerpo, pero menos aire entre objetos.
% La compactación procede de síntesis multipanel, no de reducir la legibilidad.
% --------------------------------------------------------------------------
\setlength{\parindent}{0pt}
\setlength{\parskip}{2.2pt}
\setlength{\textfloatsep}{5pt plus 1pt minus 1pt}
\setlength{\floatsep}{4pt plus 1pt minus 1pt}
\setlength{\intextsep}{4pt plus 1pt minus 1pt}
\setlength{\abovedisplayskip}{4.5pt}
\setlength{\belowdisplayskip}{4.5pt}
\setlength{\abovedisplayshortskip}{3pt}
\setlength{\belowdisplayshortskip}{3pt}
\setlength{\tabcolsep}{3.2pt}
\captionsetup{font=footnotesize,skip=2pt}

\renewcommand{\viuownsource}{%
  \par\vspace{0.4mm}\noindent
  \makebox[\linewidth][c]{\scriptsize\itshape Elaboración propia.}\par
}

% Encabezado de apendice: numeracion propia (A, B, ...) y etiqueta referenciable.
\newcommand{\anxsection}[2]{%
  \par\vspace{2mm}%
  \refstepcounter{section}%
  \noindent{\color{VIUOrange}\sffamily\bfseries\large
    Anexo~\thesection.\ #1\par}%
  \label{#2}\vspace{1mm}%
}

\usepackage{titlesec}
\newcommand{\leveltag}[1]{%
  \noindent\colorbox{VIUDark}{%
    \parbox{\dimexpr\textwidth-2\fboxsep\relax}{%
      \color{white}\sffamily\bfseries\small #1}}\par\vspace{1.2mm}
}

% La plantilla de la VIU numera los apartados y los lleva al indice. Se
% conserva el color y la tipografia que el capitulo ya usaba en sus rotulos.
\titleformat{\section}
  {\color{VIUOrange}\sffamily\bfseries\Large}{\thesection}{0.6em}{}
\titlespacing*{\section}{0pt}{2.2ex plus .6ex minus .2ex}{1.1ex plus .2ex}

% Pie de figura para material derivado de metodos publicados. La revision de
% originalidad exige que un diagrama que resume un algoritmo ajeno lo atribuya.
\newcommand{\viuderivedsource}[1]{%
  \par\vspace{0.4mm}\noindent
  \makebox[\linewidth][c]{\scriptsize\itshape Elaboración propia a partir de #1.}\par
}

\newcommand{\compactsource}{%
  {\centering\scriptsize\itshape Elaboración propia a partir de los datos experimentales.\par}
}

\newtheorem{n1prop}{Proposición}
\newtheorem{n1lema}{Lema}
\theoremstyle{remark}
\newtheorem{n1obs}{Observación}
\theoremstyle{plain}

% Panel de figura PDF. El segundo argumento es la altura MÁXIMA admitida, no una
% reserva fija: una minipage de altura fija dejaba una banda vacía cada vez que el
% asset resultaba limitado por el ancho, que es el caso de casi todas las figuras
% anchas y bajas de este capítulo.
\newcommand{\pdfpanel}[4]{%
  \begin{minipage}[t]{#1}
    \centering
    {\scriptsize\bfseries #4\par\vspace{0.5mm}}
    \IfFileExists{#3}{%
      \includegraphics[max width=\linewidth,max totalheight=#2]{#3}%
    }{%
      \fbox{\parbox[c][\dimexpr#2-8mm\relax][c]{.92\linewidth}{%
        \centering\scriptsize Asset pendiente:\\\texttt{\detokenize{#3}}}}%
    }
  \end{minipage}%
}

\newcommand{\tikzpanel}[4]{%
  \begin{minipage}[t]{#1}
    \centering
    {\scriptsize\bfseries #4\par\vspace{0.5mm}}
    % Compatibilidad con los TikZ heredados de la versión larga.
    % protocol_pipeline.tex espera \spxunitstat y n2_model_boundary.tex
    % espera \spxunitntwo definidos por el documento padre.
    \pgfmathsetlengthmacro{\spxunitstat}{\linewidth/13.5}%
    \pgfmathsetlengthmacro{\spxunitntwo}{\linewidth/13.5}%
    \IfFileExists{#3}{%
      \adjustbox{max width=\linewidth,max totalheight=#2}{\input{#3}}%
    }{%
      \fbox{\parbox[c][\dimexpr#2-8mm\relax][c]{.92\linewidth}{%
        \centering\scriptsize TikZ pendiente:\\\texttt{\detokenize{#3}}}}%
    }
  \end{minipage}%
}

% Variante para diagramas cuya fuente TikZ es menor que la caja de texto.
% \tikzpanel solo reduce; esta versión amplía de forma uniforme hasta el ancho
% disponible y evita que las etiquetas queden pequeñas en páginas dedicadas.
\newcommand{\tikzpanelfit}[4]{%
  \begin{minipage}[t]{#1}
    \centering
    {\scriptsize\bfseries #4\par\vspace{0.5mm}}
    % Mismas unidades heredadas que \tikzpanel: sin ellas, cualquier TikZ
    % antiguo que las use falla con "Undefined control sequence".
    \pgfmathsetlengthmacro{\spxunitstat}{\linewidth/13.5}%
    \pgfmathsetlengthmacro{\spxunitntwo}{\linewidth/13.5}%
    \IfFileExists{#3}{%
      \resizebox{\linewidth}{!}{\input{#3}}%
    }{%
      \fbox{\parbox[c][\dimexpr#2-8mm\relax][c]{.92\linewidth}{%
        \centering\scriptsize TikZ pendiente:\\\texttt{\detokenize{#3}}}}%
    }
  \end{minipage}%
}

\newcommand{\tightcaption}[1]{%
  \vspace{1mm}{\footnotesize #1\par}\viuownsource
}

\newcommand{\proofsketch}[1]{%
  {\small\textit{Demostración (síntesis).} #1\hfill$\square$\par}
}

% Con [H] cada figura quedaba clavada en su punto y dejaba hueco al pie.
% Estos margenes permiten que una figura grande comparta pagina con texto.
\renewcommand{\topfraction}{0.92}
\renewcommand{\bottomfraction}{0.85}
\renewcommand{\textfraction}{0.08}
\renewcommand{\floatpagefraction}{0.80}

% La cabecera oficial reserva su linea central para identificar el trabajo.
\setviuheadertitle{Trabajo Fin de M\'aster \textbar{} Jorge Luis Mayorga Taborda}

\begin{document}
\viuromanfoliofalse
\pagenumbering{arabic}
\setcounter{page}{1}

% ==========================================================================
% PÁGINA 1/20 · Problema + escalera N1--N4
% ==========================================================================
\section{SP1: formación de coaliciones de AMR}

Antes de mover una carga hay que decidir qué robots móviles autónomos (AMR) se
encargan de ella \citep{gerkeyMataric2004MRTA,korsahStentzDias2013iTax}. Cuando
todos los AMR tienen la misma capacidad, la decisión se reduce a cuántos se
reclutan, y repartir puestos intercambiables admite solución exacta
(Figura~\ref{fig:modelos}). En cuanto las capacidades difieren, el número deja de
determinar la factibilidad y pasa a importar qué AMR concretos componen la
coalición \citep{zhang2024coalition}.

SP1 estudia el reclutamiento de coaliciones de AMR heterogéneos desde cuatro
niveles sucesivos. N1 establece cuándo basta una representación por cardinalidad;
N2 conserva la capacidad individual y la atomicidad; N3 limita la información a
la comunicación vecinal; y N4 estudia cuánto orden estratégico de revisión hace
falta. Sobre ese mismo problema se evalúan además las relajaciones poblacionales
y primal-dual como alternativas continuas, en el mismo banco y con el mismo
oráculo.

La capacidad de una coalición se modela aquí como la suma de las capacidades
individuales, $Q_k=\sum_{i:a_i=k}c_i^{\mathrm{pay}}$. En transporte cooperativo esa
suma no pasa de ser una cota superior, porque la carga útil efectiva depende del
agarre y del reparto del \emph{wrench} \citep{an2023cooperativeReview}. Aun así, las dos preguntas se tratan por
separado. SP1 decide sobre el modelo aditivo y termina al verificar la coalición
entera. Comprobar si esa coalición es realizable exige modelar contacto, geometría de
agarre y reparto de \emph{wrench} ---el vector conjunto de fuerza y momento---,
por lo que se reserva para SP2; a este nivel interesa llegar allí con los
candidatos ya filtrados por una condición escalar. Los cuatro supuestos se retiran uno por
nivel.

{\small
\begin{minipage}[t]{.49\linewidth}
\textbf{Homogéneo: expansión por puestos.}
\begin{equation}
\label{eq:sp1-homogeneous-model}
\begin{aligned}
\min_x\;&\sum_{i\in\mathcal I}\sum_{h\in\mathcal H}d_{ih}x_{ih}\\
\mathrm{s.a.}\;&\sum_hx_{ih}\le1,\quad
\sum_ix_{ih}=1,\quad x_{ih}\in\{0,1\}.
\end{aligned}
\tag{1}
\end{equation}
$\mathcal H_k$ es el conjunto de puestos de la carga $k$, con
$|\mathcal H_k|=\lceil m_k/\overline{c}^{\mathrm{pay}}\rceil$ y
$c_i^{\mathrm{pay}}=\overline{c}^{\mathrm{pay}}$.
\end{minipage}\hfill
\begin{minipage}[t]{.49\linewidth}
\textbf{Heterogéneo: MILP.}
\begin{equation}
\label{eq:sp1-heterogeneous-model}
\begin{aligned}
\min_{x,e}\;&\sum_{i,k}d_{ik}x_{ik}
+\lambda_e\sum_ke_k+\lambda_n\sum_{i,k}x_{ik}\\
\mathrm{s.a.}\;&\sum_kx_{ik}\le1,\quad
e_k=\sum_ic_i^{\mathrm{pay}}x_{ik}-m_k\ge0,\\
&x_{ik}\in\{0,1\}.
\end{aligned}
\tag{2}
\end{equation}
En el diagnóstico N1.E4, $\lambda_e=0{,}25$ m/kg y
$\lambda_n=1$ mm/AMR.
\end{minipage}
}

\vspace{1mm}
\begin{figure}[htbp]\centering
\tikzpanel{.98\linewidth}{8.0cm}{figures/compact/fig01_model_boundary.tex}
{Cardinalidad frente a capacidad individual}
\caption{Expansión por puestos y formulación con capacidad individual sobre la
misma instancia: coaliciones admisibles bajo cada representación.}
\label{fig:modelos}
\viuownsource
\end{figure}

% ==========================================================================
% PÁGINA 2/20 · Protocolo común
% ==========================================================================
Cinco geometrías normalizadas y un tratamiento con retirada de un AMR generan
los mundos. Sus parámetros, la regla de inferencia por tipo de respuesta y el
protocolo Monte Carlo pareado pertenecen al capítulo de metodología. Cada semilla genera un mundo
independiente, y el mundo es la unidad de pareamiento, remuestreo e inferencia.
Dentro de un mismo mundo todos los métodos reciben los mismos AMR, las mismas
cargas y los mismos parámetros. El \emph{bootstrap}
remuestrea mundos completos por celda y Holm corrige cada familia de contrastes.
Nada se descarta, tampoco las ejecuciones fallidas o censuradas.

% PÁGINA 3/20 · N1 teoría
% ==========================================================================
\section{N1: dominio de validez de la cardinalidad}

% Página densa en teoría: se comprime el aire entre displays, no el cuerpo.
\setlength{\abovedisplayskip}{2.5pt}\setlength{\belowdisplayskip}{2.5pt}

Con capacidad común y tareas que requieren varios AMR
\citep{korsahStentzDias2013iTax}, cubrir $m_k$ equivale a reclutar
$n_k=\lceil m_k/\overline{c}^{\mathrm{pay}}\rceil$ AMR:
\begin{equation}
\label{eq:n1-cardinal}
\min_{y\in\{0,1\}^{N\times K}}\sum_{i,k}d_{ik}y_{ik},
\qquad \sum_ky_{ik}\le1,\quad \sum_iy_{ik}=n_k .
\tag{3}
\end{equation}
La expansión en $M=\sum_kn_k$ puestos intercambiables produce:
\begin{equation}
\label{eq:n1-lsap}
\min_{x\in\{0,1\}^{N\times M}}\sum_{i,h}d_{ih}x_{ih},
\qquad \sum_ix_{ih}=1,\quad \sum_hx_{ih}\le1 .
\tag{4}
\end{equation}

\begin{n1prop}[Equivalencia exacta de la expansión por puestos]
Si $c_i^{\mathrm{pay}}=\overline{c}^{\mathrm{pay}}$, los costes se replican por puesto y
$M\le N$, \eqref{eq:n1-cardinal} y \eqref{eq:n1-lsap} tienen igual factibilidad y
valor óptimo.
\end{n1prop}
\proofsketch{Una biyección entre cada coalición $S_k$ y sus puestos conserva coste;
la suma inversa $y_{ik}=\sum_{h\in\mathcal H_k}x_{ih}$ recupera la coalición.}

\begin{n1lema}[Ventana de validez de la cardinalidad]
Sea $n_k$ el número de AMR que recluta el modelo de cardinalidad, y sean
$c_{\min}$ y $c_{\max}$ la capacidad mínima y máxima de la flota candidata. Si
\begin{equation}
(n_k-1)c_{\max}<m_k\le n_kc_{\min},
\tag{5}
\end{equation}
toda coalición de $n_k$ AMR cubre $m_k$ y ninguna de $n_k-1$ puede hacerlo.
\end{n1lema}
\proofsketch{Una coalición de $n_k$ AMR aporta al menos $n_kc_{\min}\ge m_k$, luego
cubre. Una de $n_k-1$ aporta como mucho $(n_k-1)c_{\max}<m_k$, luego no puede.}
La ventana es suficiente, no necesaria. Su amplitud en $m_k$ vale
$c_{\max}-n_k(c_{\max}-c_{\min})$ y decrece linealmente con $n_k$ hasta anularse.
Con dispersión relativa $\delta=(c_{\max}-c_{\min})/c_{\max}$, es no vacía si y
solo si
\begin{equation}
\label{eq:n1-spread}
\delta<\frac{1}{n_k}.
\tag{5b}
\end{equation}

La matriz de restricciones de \eqref{eq:n1-lsap} es la incidencia de un grafo
bipartito. Al ser totalmente unimodular, su relajación lineal $x_{ih}\in[0,1]$
alcanza el óptimo en un vértice entero:
\begin{equation}
\label{eq:n1-lp-integral}
x^\star_{\rm LP}\in\{0,1\}^{N\times M},
\qquad J^{\rm puestos}_{\rm LP}=J_{\rm LSAP}.
\tag{5c}
\end{equation}
El método húngaro \citep{kuhn1955Hungarian} y las variantes de camino aumentante
\citep{jonkerVolgenant1987Shortest,crouse2016Rectangular} explotan esta estructura.
Esa estructura depende de que los puestos sean intercambiables. En cuanto la
capacidad de cada AMR entra directamente en las restricciones deja de existir, en
general, una representación por puestos equivalente, y la relajación puede
volverse fraccionaria.


\setlength{\abovedisplayskip}{4.5pt}\setlength{\belowdisplayskip}{4.5pt}

\needspace{5\baselineskip}

% ==========================================================================
% N1 · evidencia decisiva
% ==========================================================================
\section{Validez empírica de la reducción por cardinalidad}

\begin{figure}[htbp]\centering
\pdfpanel{.98\linewidth}{10.5cm}{assets/figures/sp1-final/f3_n1_validity_boundary.pdf}
{Dominio de validez de la reducción por cardinalidad}
\caption{Falsos positivos de factibilidad del modelo por puestos, por nivel de
dispersión de capacidad y geometría (N1.E4). Barras: proporción sobre
\NOneWorldsPerCell{} mundos por celda, con IC 95\,\% de Wilson. El panel inferior
recoge la severidad mediana del déficit relativo en los mundos afectados.}
\label{fig:n1}
\viuownsource
\end{figure}

Llamamos falso positivo de factibilidad a una coalición que el modelo de
cardinalidad acepta pese a incumplir la capacidad agregada real. Ninguna de las
cinco geometrías produjo uno solo mientras la capacidad fue homogénea
(Figura~\ref{fig:n1}). Al dispersarla, la proporción crece en las cinco a la vez y
llega al \NOneExtremeFalsePct\,\% en el nivel extremo. Los fallos también se
agravan. Entre el nivel bajo y el extremo, la severidad mediana del déficit pasa
del \NOneSeverityLow{} al \NOneSeverityExtreme\,\% de la demanda, tomando en
ambos casos la mediana entre las cinco geometrías. La tasa creció de forma monótona con el CV
($\beta_{\rm CV}=\NOneTrendSlope$ en escala logit, IC 95\,\%:
\NOneTrendLow--\NOneTrendHigh). El barrido de dispersión es un factor del diseño,
de modo que la relación es de asociación dentro de este banco.

Resolver la asignación por puestos de forma exacta ahorró \NOneSavingPct\,\% de
recorrido frente a asignar cada carga por cercanía (IC 95\,\%:
\NOneSavingLowPct--\NOneSavingHighPct\,\%), y pasó de cinco puntos porcentuales
---el umbral de relevancia práctica que este capítulo fija por convención--- en
tres de las cinco geometrías. La ventaja depende de con qué se compare. Al
aleatorizar el orden de recorrido de la heurística voraz, la mediana agregada baja
hasta \NOneOrderOverallPct\,\% y queda por debajo de ese mismo umbral, de modo que
el ahorro atribuible a resolver el LSAP de forma exacta no se sostiene frente a
una línea base sin orden privilegiado.


% ==========================================================================
% PÁGINA 6/20 · N2 teoría
% ==========================================================================
\section{N2: asignación con capacidades individuales}

La formulación entera conserva la identidad de cada AMR:
\begin{equation}
\label{eq:n2-milp}
\min_{y\in\{0,1\}^{N\times K}}\sum_{i,k}d_{ik}y_{ik},
\qquad \sum_ky_{ik}\le1,\quad
\sum_ic_i^{\mathrm{pay}}y_{ik}\ge m_k .
\tag{6}
\end{equation}

El objetivo de \eqref{eq:n2-milp} es la distancia pura. La formulación
\eqref{eq:sp1-heterogeneous-model} añade en cambio los términos $\lambda_e$ y
$\lambda_n$, que penalizan el exceso de capacidad y el tamaño de la coalición, y
solo se emplea como diagnóstico en N1.E4. Ahí cambian a la vez el modelo y el
criterio de optimización, así que sus valores no admiten comparación con los
de N2.

\begin{n1prop}[N1 como caso particular de N2]
Si $c_i^{\rm pay}=\overline{c}^{\rm pay}$, $d_{ik}>0$ y $\sum_kn_k\le N$, el óptimo
de \eqref{eq:n2-milp} recluta exactamente $n_k$ AMR y coincide con el LSAP.
\end{n1prop}
\proofsketch{Con capacidad común, $\sum_i\overline{c}^{\rm pay}y_{ik}\ge m_k$ equivale
sobre enteros a $\sum_iy_{ik}\ge\lceil m_k/\overline{c}^{\rm pay}\rceil=n_k$. Como
$d_{ik}>0$, todo AMR que exceda esa cota aumenta el objetivo, luego el óptimo
satisface la igualdad y \eqref{eq:n2-milp} se reduce a \eqref{eq:n1-cardinal}.}

\begin{n1obs}[Mismo entero, distinta relajación]
Relajar \eqref{eq:n2-milp} admite cubrir una carga con una fracción de AMR; la
formulación extendida por puestos sigue siendo integral.
\end{n1obs}
\proofsketch{Para una carga aislada, la cobertura relajada se satisface con
$\sum_iy_{ik}=m_k/\overline{c}^{\rm pay}$, en general fraccionario, de modo que
el óptimo del programa relajado no tiene por qué caer en un vértice entero. En
\eqref{eq:n1-lsap} la demanda ya está discretizada en puestos y la matriz es
totalmente unimodular, luego el vértice óptimo sí lo es.}
La igualdad de capacidades no basta para que la relajación sea integral. Con
$\mathrm{CV}=0$, agrupando los dos niveles de presión, la brecha mediana entre
\eqref{eq:n2-milp} y su relajación fue del \NTwoGapHomogeneousPct\,\%, porque un
AMR puede repartirse entre varias cuotas. La
integralidad de N1 procedía de la formulación extendida por puestos, y esa
formulación no admite equivalente cuando cada capacidad debe conservarse.

\begin{n1lema}[Capacidad agregada]
$\sum_i c_i^{\rm pay}\ge\sum_km_k$ es necesaria, no suficiente.
\end{n1lema}
\proofsketch{Es necesaria porque $\sum_ky_{ik}\le1$ impide que un AMR contribuya a
más de una carga, luego la capacidad total empleada no supera $\sum_ic_i^{\rm
pay}$. No es suficiente: el panel izquierdo de la
Figura~\ref{fig:n2-contraejemplos} exhibe una instancia con igualdad exacta y sin
reparto entero factible.}

\begin{n1lema}[Relajación no realizable]
Permitir $y_{ik}\in[0,1]$ puede reducir estrictamente el óptimo y volver factible
una instancia sin solución atómica.
\end{n1lema}
\proofsketch{Basta un contraejemplo de cada tipo, y el panel derecho de la
Figura~\ref{fig:n2-contraejemplos} los reúne: fraccionar un AMR entre cuotas
alcanza un coste que ninguna asignación entera reproduce, y la instancia del panel
izquierdo es factible en la relajación e infactible sobre enteros.}

La formulación ponderada ya no conserva, en general, la estructura totalmente
unimodular de N1, y su factibilidad tampoco es fácil de decidir.

\begin{n1lema}[Dureza de la factibilidad]
Decidir si \eqref{eq:n2-milp} admite solución factible es NP-difícil.
\end{n1lema}
\proofsketch{Reducción desde PARTICIÓN \citep{karp1972reducibility}. Dados
enteros $a_1,\ldots,a_N$ con $\sum_ia_i=2B$, tómense $K=2$ cargas,
$c_i^{\rm pay}=a_i$ y $m_1=m_2=B$. Como $\sum_ky_{ik}\le1$, ningún AMR
contribuye a las dos cargas, de modo que las capacidades empleadas por ambas
coaliciones son disjuntas y suman a lo sumo $2B$. Cubrirlas exige
$\sum_{i\in S_1}a_i\ge B$ y $\sum_{i\in S_2}a_i\ge B$ con
$S_1\cap S_2=\varnothing$, y con total exactamente $2B$ ambas desigualdades solo
pueden cumplirse con igualdad. Existe asignación factible si y solo si los $a_i$
admiten partición en dos subconjuntos de suma $B$. La construcción es
polinómica. Una carga sola no bastaría: asignarle toda la flota siempre
cubriría la demanda.}

El oráculo de este capítulo resuelve por tanto un problema NP-difícil, y solo
dentro del límite de tiempo declarado. De ahí que la oferta total tampoco
determine por sí sola la factibilidad cuando las capacidades difieren: dos
instancias mínimas bastan para verlo (Figura~\ref{fig:n2-contraejemplos}).

\begin{figure}[htbp]\centering
\tikzpanelfit{.98\linewidth}{8.5cm}{figures/n2_model_boundary.tex}
{Dos contraejemplos mínimos}
\caption{Instancias con capacidad agregada suficiente y sin partición entera
factible (izquierda), y con óptimo continuo inalcanzable para cualquier coalición
entera (derecha).}
\label{fig:n2-contraejemplos}
\viuownsource
\end{figure}

En la primera, la capacidad total iguala la demanda y aun así no existe reparto
entero. El LP de la segunda usa una fracción del segundo AMR y alcanza un coste
que ninguna coalición entera reproduce.

\needspace{5\baselineskip}

% ==========================================================================
% PÁGINA 7/20 · N2 E1--E2
% ==========================================================================
\section{Validación del oráculo y brecha de integralidad}

\begin{figure}[htbp]\centering
\pdfpanel{.98\linewidth}{12.5cm}{assets/figures/sp1-final/f4_n2_atomicity_map.pdf}
{Mapa de atomicidad de N2}
\caption{(A) Óptimo de HiGHS frente a enumeración exhaustiva. (B) Brecha mediana
de integralidad por CV nominal y presión $\rho$. (C) Porcentaje de mundos con LP
factible y MILP infactible. (D) Proporción de mundos con solución entera factible y
certificada por el oráculo; el descenso con la presión mezcla instancias
realmente infactibles con instancias no certificadas dentro del límite de tiempo. (E) Distribución agregada de la brecha; la línea marca la mediana.}
\label{fig:n2}
\viuownsource
\end{figure}

HiGHS \citep{huangfuHall2018Highs} hace de oráculo en todo el capítulo, y su
valor se contrastó con enumeración exhaustiva en el rango de tamaños donde esta
sigue siendo abordable. Sobre \NTwoOracleRows{} mundos de
hasta \NTwoOracleMaxStates{} asignaciones coincidió con la enumeración exhaustiva
en el \NTwoOracleAgreementPct\,\% de los casos, con
una discrepancia máxima de \NTwoOracleMaxErrorPm{} pm, que es redondeo en coma
flotante y no una diferencia de recorrido (Figura~\ref{fig:n2}A).

El LP puede asignar una fracción de un AMR, mientras que una coalición ejecutable
requiere que el AMR entre entero o no entre. Sobre \NTwoAtomicityRows{} mundos el
óptimo continuo
resultó fraccionario en el \NTwoFractionalAllPct\,\%, y en los
\NTwoAtomicityComparable{} pares resolubles la brecha nunca fue nula. Su mediana
alcanzó \NTwoGapPct\,\% (IC 95\,\%:
\NTwoGapLowPct--\NTwoGapHighPct\,\%).

La brecha no es monótona en la heterogeneidad. Es máxima en los dos extremos del
barrido y mínima en el centro: con $\rho=0{,}70$ vale \NTwoGapCvZeroRho{} con
$\mathrm{CV}=0$, baja a \NTwoGapCvLowRho{} y \NTwoGapCvMidRho, repunta a
\NTwoGapCvHighRho{} con $\mathrm{CV}=0{,}65$ y vuelve a \NTwoGapCvOneRho\,\% con
$\mathrm{CV}=1$.

Los \NTwoLpOnlyCount{} mundos que el LP resuelve y el MILP no se reparten de forma
muy desigual entre las celdas del barrido. \NTwoLpOnlyTopN{} caen en la celda homogénea de mayor presión,
$\mathrm{CV}=0$ con $\rho=0{,}85$, y los \NTwoLpOnlyRestN{} restantes en
$\mathrm{CV}=1$ con esa misma presión. La heterogeneidad no hace falta para que LP
y MILP difieran.

Hasta dónde llega el referente lo marca un barrido de presión aparte, con
\NOneWorldsPerCell{} mundos por celda, en el que la proporción con solución entera
factible y certificada es completa hasta $\rho=0{,}75$, cae al
\NTwoCertHomogeneousHighRho\,\% en la celda homogénea de $\rho=0{,}85$ y se anula
a partir de $\rho=0{,}92$ con $\mathrm{CV}=0$. Ese descenso mezcla dos cosas
distintas: instancias realmente infactibles e instancias que el solver no logró
certificar dentro del límite de tiempo. Las brechas de este capítulo se calculan
solo donde el oráculo certificó.

Hasta aquí el optimizador ha conocido las capacidades y posiciones de toda la
flota sin pagar por reunirlas. N3 retira ese supuesto y cobra cada byte que cruza
una arista. La pregunta pasa de qué coalición es óptima a qué coalición puede
acordarse con información vecinal.

% ==========================================================================
% PÁGINA 9/20 · N3 contrato y frontera teórica
% ==========================================================================
\section{N3: reclutamiento con información vecinal y métodos comparados}

Cada AMR conoce su capacidad y su posición. El catálogo de cargas es público. Todo
lo demás ---el estado de cualquier otro AMR--- tiene que cruzar una arista, y cada
byte que cruza se contabiliza. El observador se limita a registrar la ejecución y
a calcular las métricas cuando termina.

El grafo de comunicación es de disco. Dos AMR son vecinos cuando
$d_{ij}\le\min(R_i^{\rm com},R_j^{\rm com})$, y $\mathcal N_i$ denota los vecinos
de $i$. La formulación admite radios de comunicación distintos por AMR; esa
generalización se reserva para SP4. En SP1 la flota entera comparte
$R=m\,r_c$. La escala $r_c$ es la mayor arista del
árbol generador mínimo euclídeo del despliegue, es decir, el radio exacto al que el
grafo pasa a ser conexo. Los cuatro multiplicadores
$m\in\{0{,}95;\,1{,}05;\,1{,}50;\,2{,}00\}$ caen a uno y otro lado de esa
transición, y el grafo completo cierra la serie como quinto régimen.


Dos agregados bastan entonces para decidir si una asignación es factible. Uno mide
el déficit de capacidad. El otro, el exceso de compromiso.

\begin{n1prop}[Certificado de factibilidad]
Sea $y\in\{0,1\}^{N\times K}$ una asignación candidata cualquiera, no
necesariamente factible para \eqref{eq:n2-milp}. Con
$D(y)=\sum_k[m_k-Q_k(y)]_+$ y $C(y)=\sum_i[\sum_ky_{ik}-1]_+$, esa asignación es
factible si y solo si $D(y)=C(y)=0$.
\end{n1prop}
\proofsketch{Ambos términos son sumas de partes positivas, luego se anulan si y
solo si se anula cada sumando. $D(y)=0$ equivale a $Q_k(y)\ge m_k$ para toda carga,
y $C(y)=0$ equivale a $\sum_ky_{ik}\le1$ para todo AMR. Esas dos familias son
exactamente las restricciones de \eqref{eq:n2-milp}.}
El exceso $C$ solo interviene mientras los métodos de N3 no han acordado un
perfil; en N4 la variable $a_i$ lo anula por construcción. Ese certificado se
evalúa con lo que cada AMR sabe, y ahí aparece un límite que no depende del
método: bajo partición permanente ninguna política que use solo la información
de su componente puede garantizar el óptimo global. El argumento es de
indistinguibilidad, en el sentido de \citet{halpernMoses1990Knowledge}.

\begin{n1lema}[Indistinguibilidad local]
Con un grafo permanentemente desconectado, ninguna política determinista basada
solo en la vista de su componente garantiza el óptimo global en todas las
instancias factibles.
\end{n1lema}
\proofsketch{Sean $W$ y $W'$ dos mundos que coinciden en el catálogo público de
cargas y en todo lo que ocurre dentro de la componente $C_1$, y que difieren
únicamente en las capacidades privadas $c_j^{\rm pay}$ de los AMR de $C_2$: en $W$
esas capacidades cubren la demanda asignada a $C_2$ y en $W'$ no. Bajo partición
permanente ningún mensaje cruza entre componentes, de modo que la vista de todo
AMR de $C_1$ ---su propio estado, sus aristas y el catálogo--- es idéntica en
ambos mundos, y una política determinista elige la misma acción. El óptimo social,
en cambio, exige en $W'$ que $C_1$ ceda capacidad a las cargas que $C_2$ no puede
cubrir, y en $W$ no: como $d_{ik}>0$, enviar esa capacidad en $W$ es innecesario
y aumenta estrictamente el coste, mientras que en $W'$ es la única forma de
alcanzar la factibilidad. Una misma acción no puede ser óptima en ambos mundos,
de modo que la política falla en al menos uno.}

La campaña de N3 evaluó \NThreeWorlds{} mundos con $N=16$ y $K=5$, de los que el
oráculo certificó \NThreeOracleCertified{} y demostró la infactibilidad de
\NThreeOracleInfeasible. Los tres referentes que se comparan en ellos proceden de
dos familias distintas. De la subasta con consenso de CBBA
\citep{choiBrunetHow2009CBBA,johnson2011acbba} sale CBBA-RB, que añade una barrera
de retorno; Weighted-GRAPE y Pair-GRAPE vienen en
cambio de la formación hedónica de GRAPE \citep{jangShinTsourdos2018GRAPE,dutta2021hedonic},
cuya noción de estabilidad procede de \citet{bogomolnaiaJackson2002Hedonic}.
Las tres variantes comparten
la cobertura ponderada por capacidad. Su pseudocódigo se recoge en el
Anexo~\ref{anx:pseudocodigo}.


{\footnotesize La ventana sin cambios de CBBA-RB es un criterio de parada, no una
garantía; las del CBBA original no se trasladan a esta extensión multi-ganador.
Weighted-GRAPE sí la tiene con entrega fiable, porque cada perfil aceptado reduce
estrictamente $(D,J)$.}

\needspace{5\baselineskip}

% ==========================================================================
% PÁGINA 11/20 · N3 resultados
% ==========================================================================
\section{Calidad, comunicación y conectividad}

\begin{figure}[htbp]\centering
\pdfpanel{.96\linewidth}{7.2cm}{assets/figures/sp1-final/f5_n3a_quality_communication.pdf}
{Calidad y coste de comunicación}
\caption{Factibilidad condicionada al oráculo sobre \NThreeWorlds{} mundos,
brecha mediana sobre los \SPoneNThreeCommonWorlds{} de soporte común y bytes/AMR
de los tres métodos (N3.E2). Intervalos de Wilson para la factibilidad y
\emph{bootstrap} pareado para la brecha.}
\label{fig:n3-resultados}
\viuownsource
\end{figure}

CBBA-RB mantiene una tabla parcial de ganadores. Las dos variantes GRAPE, que
reconstruyen un perfil común completo, fallaron la cobertura menos que CBBA-RB (Figura~\ref{fig:n3-resultados}) y resultaron factibles en prácticamente
todos los mundos con solución ---\NThreeFeasPctGrape{} y \NThreeFeasPctPair\,\%---
frente al \NThreeFeasPctCbba\,\% de CBBA-RB. Las brechas medianas sobre los
\SPoneNThreeCommonWorlds{} mundos en que los tres entregan salida factible fueron
\SPoneNThreeCbbaGap, \SPoneNThreeGrapeGap{} y \SPoneNThreePairGrapeGap\,\%. Ese
orden se mantuvo mayoritariamente de un mundo a otro ($W=\NThreeKendallW$).

CBBA-RB emitió \NThreeBytesCbba{} bytes/AMR; Weighted-GRAPE y Pair-GRAPE,
\NThreeBytesGrape{} y \NThreeBytesPair.

El barrido de topología es una campaña distinta de la anterior, con
\NThreeTopoWorldsPerRegime{} mundos por régimen, de modo que sus factibilidades no
son comparables término a término con las de la Figura~\ref{fig:n3-resultados}.
Dentro de ella, más conectividad no dio mejor cobertura a CBBA-RB: su factibilidad decrece de forma monótona desde el
\NThreeTopoCbbaThreshold\,\% cerca del umbral hasta el
\NThreeTopoCbbaComplete\,\% con el grafo completo, pasando por
\NThreeTopoCbbaMedium{} y \NThreeTopoCbbaDense\,\%. Las dos variantes GRAPE se
mantuvieron entre \NThreeTopoGrapeMin{} y \NThreeTopoGrapeMax\,\% en los cuatro
regímenes conexos. El tráfico decreció al ralear el grafo en los tres métodos.

Bajo partición permanente la cobertura cae al \NThreeTopoCbbaPartitioned\,\% en
CBBA-RB, pero ese régimen es el control negativo y su barra mide cobertura
verificable \emph{post hoc}. Las variantes GRAPE terminaron ahí con perfiles
locales distintos entre componentes. Ese último perfil puede medirse, y no
constituye una coalición acordada.

Como diagnóstico posterior se comprobó si alguna medida agregada ordena los
regímenes mundo a mundo. Ni el índice adimensional de los grafos geométricos
aleatorios \citep{penrose2003RandomGeometric} ni la conectividad algebraica
$\lambda_2$ lo consiguen: los rangos de dos regímenes contiguos se solapan en
ambos casos, y el desarrollo se recoge en el Anexo~\ref{anx:conectividad}. La
etiqueta de régimen describe entonces una consigna de diseño experimental y no
mide la conectividad que cada mundo acaba teniendo.



% ==========================================================================
% Comparadores continuos: formulacion
% El encabezado abre pagina: si no, queda solo al pie de la pagina anterior,
% separado del desarrollo.
% ==========================================================================
\needspace{5\baselineskip}
\section{Dinámicas continuas, cierre entero y orden de mérito}

Las dinámicas poblacionales entran en SP1 por dos vías. La primera las usa como
relajación continua: cada AMR reparte una unidad de masa sobre las cargas y
sigue una de cuatro dinámicas ---replicador, Smith, BNN y Logit--- sobre un
potencial que penaliza el déficit de cobertura y la distancia
\citep{sandholm2010population,quijano2017population}. La segunda las lleva
directamente a la decisión entera, y se trata en N4: Geo-ASR ejecuta el
protocolo de revisión de Smith en forma atómica, sin pasar por ningún estado
continuo. Junto a ellas se evalúa un esquema primal-dual en el que la cobertura
es una restricción compartida con precios sombra, cuya variante distribuida
añade un término de consenso que lleva a todos los AMR al mismo precio
\citep{yi2019operatorGNE,franci2022stochasticGNE}. Las dos formulaciones
continuas se recogen en el Anexo~\ref{anx:formulacion}.


El replicador es imitativo. No puede dar masa a una cuota con $x_{ik}=0$, de modo
que sus puntos de reposo incluyen los vértices del símplex. Smith, BNN y Logit sí
exploran cuotas sin uso previo. Ambas familias pertenecen a clases distintas de
protocolo de revisión, y sus versiones distribuidas en tiempo discreto se recogen
en \citet{martinezPiazuelo2022Population}.

Ninguna de las dos produce una coalición ejecutable: entregan participaciones
$x_{ik}\in[0,1]$. El cierre entero
$\mathcal R$ las convierte en una asignación binaria y es el mismo para todos los
métodos continuos, de modo que el posproceso se mantiene fijo entre
comparaciones. Eso no implica que les afecte por igual: su efecto depende del
perfil continuo que recibe. $\mathcal R$ opera en dos pasos. Primero
asigna cada AMR a la carga con mayor masa, y lo deja inactivo si la masa restante
$1-\sum_kx_{ik}$ supera a esa máxima. Después repara el resultado con cadenas de
reasignación acotadas: un paso intermedio puede aumentar el déficit al retirar un
AMR útil de otra coalición, y la cadena se acepta solo si el déficit normalizado
total baja de forma estricta. Las cadenas se truncan en ocho movimientos y la
exploración en 5\,000 nodos por reparación, así que $\mathcal R$ es una heurística
determinista y no un solver entero: no garantiza factibilidad. Se ejecuta de forma
central sobre el estado completo, y su tráfico queda fuera de las métricas de
comunicación.

En el esquema primal-dual, \eqref{eq:n4-primal-dual} emplea un multiplicador
central y \eqref{eq:n4-consensus} es su variante distribuida. Ahí el término
$-\gamma\sum_{j\in\mathcal N_i}(\lambda_{ik}-\lambda_{jk})$ arrastra a todos los
AMR hacia el mismo precio sombra de la cobertura. Para el juego convexo considerado y bajo condiciones de regularidad de tipo
Slater, un equilibrio de Nash generalizado variacional satisface las condiciones
de Karush--Kuhn--Tucker \eqref{eq:n4-kkt} con un multiplicador $\lambda^\star$
común a la restricción compartida
\citep{yi2019operatorGNE,franci2022stochasticGNE}. Esas
condiciones no se verifican aquí instancia a instancia. Como cada $J_i$ depende
solo de $x_i$ y el acoplamiento es la restricción común, el pseudogradiente resulta
fuertemente monótono con módulo $\epsilon$. Las dos familias cierran igual, con
$\mathcal R$, que es central y cuya comunicación no se contabiliza; el
Anexo~\ref{anx:formulacion} recoge la formulación de ambas.

El estado continuo no anticipa el resultado entero. El método con menor déficit continuo no
fue el de menor brecha tras $\mathcal R$ en el \NFourFiiMismatchPct\,\% de los
\NFourContWorldsN{} mundos. Y sobre los \NFourContSupportN{} que el oráculo
certificó como factibles las cuatro dinámicas terminaron muy juntas, entre
\NFourFiiRateMin{} y \NFourFiiRateMax\,\% de tasa entera.

El esquema primal-dual cruzó el criterio conjunto de convergencia en el
\NFourPdConvergencePct\,\% de los mundos. Tras el cierre central, su salida fue
factible en \NFourPdFeasibleN{} de esos \NFourPdSupportN{}
(\NFourPdFeasiblePct\,\%), con brecha mediana del \NFourPdClosedGapPct\,\%.
Ninguna de las dos familias entrega, por tanto, una coalición sin pasar por
$\mathcal R$.

% ==========================================================================
% PÁGINA 12/20 · N4 problema, taxonomía y atlas
% ==========================================================================
\section{N4: refinamiento coalicional de la asignación}

Weighted-GRAPE termina donde ningún AMR puede mejorar por su cuenta. Pair-GRAPE
añade una fase de desviaciones conjuntas con un vecino y obtuvo una brecha mediana
de \NThreeCommonGapPctPair\,\% frente al \NThreeCommonGapPctGrape\,\% de
Weighted-GRAPE sobre el mismo soporte. La mejora aparece mezclada con el resto del protocolo GRAPE. Para
aislarla hay que mantener fijo el criterio de mejora y variar cuántos AMR pueden
cambiar de asignación de forma conjunta.

La asignación entera $a$ debe cubrir toda carga activa:
\begin{equation}
\label{eq:n4-common-problem}
\begin{gathered}
a_i\in\{0,1,\ldots,K\},\quad
Q_k(a)=\sum_i c_i^{\rm pay}\mathbf1\{a_i=k\},\\
D_4(a)=\sum_{k\in\mathcal K_{\rm act}}[m_k-Q_k(a)]_+,\quad
J_4(a)=\sum_{i:a_i\ne0}d_{i,a_i},\\
D_4(a)=0,\qquad J_4(a)\to\min .
\end{gathered}
\tag{14}
\end{equation}
$\mathcal K_{\rm act}$ es el conjunto de cargas activas y $a_i=0$ significa no
participar. Geo-QPG decide directamente sobre $a_i$, la misma variable que recibe
SP2.

Seis referentes vienen de la literatura y aquí solo se adaptan a cobertura
ponderada: la asignación exacta por puestos
\citep{kuhn1955Hungarian}, la subasta con consenso
\citep{choiBrunetHow2009CBBA}, la formación hedónica
\citep{jangShinTsourdos2018GRAPE}, las dinámicas poblacionales
\citep{sandholm2010population}, la búsqueda de equilibrios generalizados
\citep{yi2019operatorGNE} y la selección distribuida de conjuntos independientes
\citep{luby1986MIS}. La formulación Geo-QPG sí se construye aquí. La componen el
criterio lexicográfico $\Psi_4$ de \eqref{eq:n4-lex} junto con su regla de
aceptación, la jerarquía $\mathcal V_h$ de \eqref{eq:n4-neighbourhood}, el
diagnóstico $h_c^\star$, la composición de propuestas disjuntas y una confirmación
vecinal apoyada en DMIS y \emph{commit} versionado.

\section{Revisión unilateral, revisión coalicional y concurrencia}

Las tres reglas de revisión eligen de forma distinta dentro de la misma vecindad
$\mathcal V_1(a)$: BR toma la mejor respuesta, Geo-ASR aplica la revisión de
Smith en su forma atómica ---probabilidad proporcional a la ganancia positiva--- y
Geo-LLL usa el aprendizaje log-lineal \citep{mardenShamma2012LogLinear}, que
admite además estados transitorios peores. Ninguna sale de las desviaciones de un
solo AMR, así que una barrera bilateral les resulta inaccesible. Al ampliar $h$
cambia qué grupos pueden escapar de un terminal unilateral. El
Anexo~\ref{anx:reglas} despliega las tres reglas y las dos ampliaciones sobre un
mismo estado.

DMIS+TX actúa después de generar las propuestas. Selecciona cuáles pueden
confirmarse a la vez, sin orden global, y para ello reutiliza la selección
distribuida de conjuntos independientes maximales
\citep{luby1986MIS}. Aquí \emph{commit} atómico designa una confirmación de escritura todo-o-nada
\citep{grayLamport2006Commit}, distinta de la atomicidad de la asignación, que se
refiere a la indivisibilidad del AMR. DMIS abrevia la selección distribuida de un
conjunto independiente maximal, y TX la transacción versionada que la acompaña.

% ==========================================================================
% PÁGINA 13/20 · N4 F-I potencial, vecindad y revisión
% ==========================================================================
\section{Geo-QPG: mejora marginal y jerarquía $h$-local}

Geo-QPG compara lexicográficamente:
\begin{equation}
\label{eq:n4-lex}
\Psi_4(a)=\bigl(D_4(a),J_4(a)\bigr).
\tag{15}
\end{equation}
Para una desviación unilateral $i:r\to s$:
\begin{equation}
\begin{aligned}
\Delta_iD={}&[m_r-(Q_r-c_i^{\rm pay})]_+
+[m_s-(Q_s+c_i^{\rm pay})]_+\\[-1mm]
&-[m_r-Q_r]_+-[m_s-Q_s]_+,\qquad
\Delta_iJ=d_{is}-d_{ir}.
\end{aligned}
\tag{16}
\end{equation}
Definiendo $\Phi_Q=-D_4$ y $\Phi_G=-J_4$:
\begin{equation}
\label{eq:n4-potential}
u_i^X(a)=\Phi_X(a)-\Phi_X(0,a_{-i}),\qquad X\in\{Q,G\}.
\tag{17}
\end{equation}
La utilidad de \eqref{eq:n4-potential} satisface la identidad de potencial exacto
de \citet{monderer1996potential} sobre cuotas heterogéneas.
\begin{n1prop}[Potencial exacto unilateral]
Dentro de una misma fase, la diferencia de utilidad coincide exactamente con la
diferencia de $\Phi_X$; para $h>1$ se habla de búsqueda coalicional finita.
\end{n1prop}
\proofsketch{Por \eqref{eq:n4-potential},
$u_i^X(s,a_{-i})-u_i^X(r,a_{-i})=\Phi_X(s,a_{-i})-\Phi_X(r,a_{-i})$: el término
$\Phi_X(0,a_{-i})$ no depende de la acción de $i$ y se cancela. La identidad es
algebraica y no usa homogeneidad de capacidades.}

El orden $h$ limita cuántos AMR pueden cambiar de asignación en una misma
revisión: con $h=1$ solo caben cambios individuales, con $h=2$ también los pares
y con $h=3$ las ternas. La vecindad estratégica es
\begin{equation}
\label{eq:n4-neighbourhood}
\mathcal V_h(a)=\{a':1\le|\{i:a_i'\ne a_i\}|\le h\},
\qquad
\mathcal V_1\subseteq\mathcal V_2\subseteq\mathcal V_3 .
\tag{18}
\end{equation}
Evaluar una desviación solo exige los agregados de las cargas implicadas
(Figura~\ref{fig:desviacion}). Con
$g_i^X(k)=u_i^X(k,a_{-i})-u_i^X(a_i,a_{-i})$ la ganancia de que $i$ pase a la carga
$k$, las reglas unilaterales ejecutadas fueron, con LLL en la forma
log-lineal habitual \citep{mardenShamma2012LogLinear} y $\tau_{\rm LLL}$ su
temperatura, distinta del peso $\beta$ de \eqref{eq:n4-fii-potential} y de la
pendiente $\beta_{\rm CV}$ de N1,
\begin{equation}
\begin{aligned}
\mathrm{BR}:&\quad k^\star\in\arg\max_k g_i^X(k),\ g_i^X(k^\star)>0,\\
\mathrm{ASR}:&\quad
\Pr(k)=\frac{[g_i^X(k)]_+}{\sum_q[g_i^X(q)]_+},\\
\mathrm{LLL}:&\quad
\Pr(k)=\frac{e^{u_i^X(k,a_{-i})/\tau_{\rm LLL}}}
{\sum_qe^{u_i^X(q,a_{-i})/\tau_{\rm LLL}}}.
\end{aligned}
\tag{19}
\end{equation}
Sea cual sea la regla, una desviación se confirma si y solo si reduce $\Psi_4$ en
orden lexicográfico: nunca aumenta el déficit y, a igualdad de déficit, reduce
$J_4$. Sobre esa condición se apoya la terminación del
Algoritmo~\ref{alg:refinamiento}.

\begin{figure}[htbp]\centering
\tikzpanelfit{.98\linewidth}{6.8cm}{figures/n4/n4_marginal_move.tex}
{Agregados que intervienen en una desviación}
\caption{Agregados $Q_k$ y $d_{ik}$ que intervienen en $\Delta_iD$ y
$\Delta_iJ$ para una desviación $i:r\to s$. Las vecindades $\mathcal V_1$,
$\mathcal V_2$ y $\mathcal V_3$ frente a las tres reglas se despliegan en el
Anexo~\ref{anx:reglas}.}
\label{fig:desviacion}
\viuownsource
\end{figure}

\needspace{5\baselineskip}

% ==========================================================================
% PÁGINA 14/20 · N4 barreras, refinamiento y commit
% ==========================================================================
\section{Barrera bilateral, refinamiento y terminación}

Interesa además el menor orden en el que aparece una mejora cuyos AMR implicados
siguen conectados entre sí, porque solo esas pueden negociarse sin intermediarios.
El orden mínimo general y su versión conectada son:
\begin{equation}
h^\star(a)=\min\{h:\exists a'\in\mathcal V_h(a),\,
\Psi_4(a')<_{\rm lex}\Psi_4(a)\}.
\tag{20}
\end{equation}
\begin{equation}
\begin{aligned}
\mathcal V_h^{\rm c}(a)&=\{a'\in\mathcal V_h(a):
G[\{i:a_i'\ne a_i\}]\ {\rm conexo}\},\\
h_c^\star(a)&=\min\{h:\exists a'\in\mathcal V_h^{\rm c}(a),\
\Psi_4(a')<_{\rm lex}\Psi_4(a)\}.
\end{aligned}
\tag{21}
\end{equation}
Una barrera bilateral suficiente satisface:
\begin{equation}
\label{eq:n4-bilateral-barrier-compact}
\begin{gathered}
Q_r-c_i^{\rm pay}<m_r,\quad Q_s-c_j^{\rm pay}<m_s,\\
Q_r-c_i^{\rm pay}+c_j^{\rm pay}\ge m_r,\quad
Q_s-c_j^{\rm pay}+c_i^{\rm pay}\ge m_s,\\
d_{is}+d_{jr}<d_{ir}+d_{js}.
\end{gathered}
\tag{22}
\end{equation}

\begin{n1lema}[Barrera unilateral y escape bilateral]
Si se cumple \eqref{eq:n4-bilateral-barrier-compact}, ninguno de los dos
movimientos que componen el intercambio mejora por separado ---cada uno crea
déficit--- y aplicarlos a la vez mantiene la cobertura y reduce $J_4$.
\end{n1lema}
\proofsketch{Las dos primeras desigualdades dan $Q_r-c_i^{\rm pay}<m_r$ y
$Q_s-c_j^{\rm pay}<m_s$: retirar $i$ de $r$, o $j$ de $s$, deja la carga
descubierta y aumenta $D_4$. Las dos siguientes garantizan que tras el intercambio
ambas quedan cubiertas, luego $\Delta D_4=0$. La última da
$\Delta J_4=(d_{is}+d_{jr})-(d_{ir}+d_{js})<0$.}

\begin{algorithm}[H]
\caption{Refinamiento progresivo por orden de vecindad.}
\label{alg:refinamiento}
\begin{algorithmic}[1]
\For{$h=1,\ldots,h_{\max}$}
\State aceptar mejoras estrictas de $\mathcal V_h$ hasta agotamiento o límite
\EndFor
\State \Return $a^{(h_{\max})}$
\end{algorithmic}
\end{algorithm}

\begin{n1prop}[Terminación sin truncamiento]
\label{prop:terminacion}
Si todos los candidatos se examinan y cada \emph{commit} reduce $\Psi_4$, el
espacio finito termina en un mínimo $h$-local.
\end{n1prop}
\proofsketch{El espacio $\{0,\ldots,K\}^N$ es finito y $<_{\rm lex}$ es un orden
estricto sobre $\Psi_4$, de modo que ningún estado puede repetirse y no hay ciclos.
La cadena termina, y al terminar ninguna desviación de $\mathcal V_h$ reduce
$\Psi_4$: eso es un mínimo $h$-local.}

\begin{n1prop}[Composición]
Propuestas que no comparten AMR ni cargas ---ni de origen ni de destino--- suman
sus variaciones; si todas mejoran, el lote también reduce $\Psi_4$.
\end{n1prop}
\proofsketch{$D_4$ suma sobre cargas y $J_4$ sobre AMR; con soportes disjuntos cada
propuesta altera sumandos distintos, luego $\Delta D_4=\sum_p\Delta D_p$ y
$\Delta J_4=\sum_p\Delta J_p$. Si toda propuesta mejora en orden lexicográfico
entonces $\Delta D_p\le0$ para todas: si alguna es negativa la suma lo es, y si
todas son nulas entonces todas tienen $\Delta J_p<0$ y la suma también.}

DMIS+TX elimina el orden global de confirmación y deja intacto el resto de la
infraestructura lógica. El grafo de conflictos
enlaza las propuestas que comparten AMR o carga, y el conjunto independiente
selecciona las que pueden confirmarse a la vez. Ninguna de las dos vías entrega
un certificado de optimalidad global.

\needspace{5\baselineskip}

% ==========================================================================
% PÁGINA 17/20 · F-I resultados E4/E9
% ==========================================================================
\section{Efecto del orden estratégico y primer escape conectado}

\begin{figure}[htbp]\centering
\pdfpanel{.98\linewidth}{12.0cm}{assets/figures/sp1-final/f7_n4_strategic_order.pdf}
{Orden estratégico y primer escape}
\caption{(A) Brecha mediana frente al MILP sobre los \SPoneCommonWorlds{} mundos
de soporte común. (B) Riesgo residual de infactibilidad sobre los
\SPoneCertWorlds{} certificados, IC de Wilson. (C) Distribución de $h_c^\star$
sobre \SPoneHstarWorlds{} terminales BR. (D) Proporción de mundos sin escape
conectado detectado hasta $h=3$: CV y presión proceden del bloque confirmatorio,
que fija $N/K=\NFourMainBlockNk$; la relación $N/K$ procede de un bloque de
sensibilidad independiente de \NFourRatioBlockN{} mundos, de modo que las tres
agrupaciones no comparten denominador. Cada barra rotula su $n$.}
\label{fig:orden-brecha}
\viuownsource
\end{figure}

Permitir que dos AMR cambien a la vez bajó la brecha mediana de
\SPoneBRGapCommon{} a \SPoneTwoBRGapCommon\,\%. Con tres quedó en
\SPoneCThreeGapCommon\,\% (Figura~\ref{fig:orden-brecha}). Las tres cifras se miden
sobre los \SPoneCommonWorlds{} mundos en que BR, 2BR y C3 entregan salida factible
y el oráculo certificó el óptimo. Entre los extremos la diferencia pareada fue de
\SPoneCThreeMinusBRCommon{} puntos (IC 95\,\%:
\SPoneCThreeMinusBRCommonLo{} a \SPoneCThreeMinusBRCommonHi).

Aparte queda la factibilidad, medida sobre los \SPoneCertWorlds{} mundos que el
oráculo declaró factibles y certificó. BR resolvió \SPoneBRFeasN{} de
\SPoneCertWorlds{} (\SPoneBRFeasRound\,\%), 2BR \SPoneTwoBRFeasN{}
(\SPoneTwoBRFeasRound\,\%) y C3 los \SPoneCertWorlds{} (100\,\%). Hasta este punto el efecto
del orden sigue confundido con el de la regla de selección, porque ambas cosas
cambiaron a la vez. La campaña factorial de la sección siguiente los separa.

Tras agotar BR, la primera mejora conectada fue bilateral en \NFourHStarTwoN{} de
los \NFourHStarWorldsN{} mundos y trilateral en \NFourHStarThreeN; en los
\NFourHStarGtThreeN{} restantes la búsqueda terminó en $h=3$ sin encontrar ninguna
(Figura~\ref{fig:orden-brecha}C). Esos \NFourHStarGtThreeN{} casos se concentran
en presión alta y capacidad homogénea, y su tasa baja al aumentar el número de AMR
por carga. El desglose por CV, presión y relación $N/K$ procede de dos bloques con
diseños distintos que no comparten denominador, y se recoge en el
Anexo~\ref{anx:sensibilidad}.

Los experimentos posteriores emplean $h=2$. La elección no se sigue de
$h_c^\star$, que localiza el primer escape y no la mejora acumulada: es un
compromiso entre las dos magnitudes medidas, ya que pasar de $h=2$ a $h=3$ reduce
la brecha mediana de \SPoneTwoBRGapCommon{} a \SPoneCThreeGapCommon\,\% y
multiplica por más de cuatro los bytes por AMR.

\needspace{5\baselineskip}

% ==========================================================================
% N4 · mecanismo y arquitectura
% ==========================================================================
\section{Magnitud de la mejora y coste de la confirmación vecinal}

$h_c^\star$ señala a qué orden aparece el primer escape y no cuánto mejora. En los
\NFourHStarLinkThreeN{} terminales con $h_c^\star=3$ la mejora adicional apareció
al habilitar desviaciones trilaterales, pero son solo \NFourHStarLinkThreeN{} de
\NFourHStarWorldsN{} mundos y de esa categoría no se extrae ninguna conclusión
estadística. El desglose por categoría, sobre los \NFourHStarLinkN{} mundos en que
BR terminó con salida factible y las tres brechas quedan definidas, se recoge en
el Anexo~\ref{anx:hstar}.

Para separar la arquitectura de confirmación del orden estratégico, la
ablación de la Figura~\ref{fig:ablacion} mantiene fijos el conjunto de candidatos
y la regla de revisión, y solo cambia cómo se confirman las propuestas: con un
orden global de confirmación (CF) o con arbitraje entre vecinos (DMIS+TX). Sus
brechas no coinciden con las de la Figura~\ref{fig:orden-brecha} porque el soporte
es el de los \SPoneDmisWorldsCert{} mundos certificados y no el soporte común de
los tres órdenes.

\begin{figure}[htbp]\centering
\pdfpanel{.94\linewidth}{7.5cm}{assets/figures/sp1-final/f8_n4_distributed_execution.pdf}
{Confirmación global frente a DMIS+TX}
\caption{Brecha, factibilidad y bytes/AMR de la confirmación con orden global y de
DMIS+TX sobre los mismos \SPoneDmisWorldsCert{} mundos certificados, y diferencia
pareada en bytes/AMR con IC 95\,\% \emph{bootstrap}.}
\label{fig:ablacion}
\viuownsource
\end{figure}

Las dos arquitecturas coinciden en los \SPoneDmisInvariantsK{} agregados que el
registro conserva de cada ejecución ---coste de distancia, déficit de capacidad,
exceso de capacidad, AMR asignados, tamaño máximo de coalición, cargas sin servir
y factibilidad--- en los \SPoneDmisSameObjN{} de \SPoneDmisPairedN{} mundos
emparejados, con una discrepancia máxima en distancia por debajo del picómetro, es
decir, redondeo en coma flotante. El registro no guarda
el vector de asignación, así que la identidad de las coaliciones no queda
demostrada; lo que sí queda es que ninguna magnitud medida las distingue. La
diferencia entre ambas está en el coste de acordarlas:
\SPoneDmisBytesDeltaCert{} bytes/AMR adicionales sobre los
\SPoneDmisWorldsCert{} certificados (IC 95\,\%:
\SPoneDmisBytesDeltaCertLo--\SPoneDmisBytesDeltaCertHi).

\section{Control del orden de revisión y banco industrial}

Se ejecutaron dos campañas adicionales con semillas no solapadas respecto de las
campañas anteriores, verificadas por resumen criptográfico del contenido de cada
mundo antes de analizar nada. La primera cruza la regla de revisión con la
función de mérito; la segunda evalúa los métodos ya publicados en cuatro
geometrías industriales que no se usaron durante el desarrollo.

El factorial emplea tres funciones de mérito: $F_0$ ordena los candidatos solo
por distancia, $F_1$ solo por cobertura y $F_2$ es la utilidad marginal de
\eqref{eq:n4-potential}. La regla de aceptación es idéntica en todas las celdas,
de modo que lo único que varía es cómo se ordenan las desviaciones ya
admisibles. A las tres reglas unilaterales se añade 2BR con $h=2$ como brazo de
control, ejecutado sobre los mismos mundos, con las mismas semillas y el mismo
presupuesto: sin él la comparación sería contra otra campaña y no admitiría
pareamiento. Son \NThreeBWorlds{} mundos y \NThreeBRuns{} ejecuciones.

Las \NThreeBUnilateralCells{} celdas unilaterales quedan entre
\NThreeBBestGap{} y \NThreeBWorstGap\,\% de brecha mediana, y las tres celdas de
control entre \NThreeBCtrlBestGap{} y \NThreeBCtrlWorstGap\,\%
(Figura~\ref{fig:n3b}A). La peor celda con $h=2$ mejora a la mejor con $h=1$ en
\NThreeBSeparation{} puntos, y la separación no se paga en cobertura: el control
también tiene la factibilidad más alta del cuadro (Figura~\ref{fig:n3b}B). El
contraste pareado por mundo confirma el efecto para cada función de mérito
(Figura~\ref{fig:n3b}C): la brecha baja \NThreeBPairedFZeroDrop{} puntos con
$F_0$ (IC 95\,\%: \NThreeBPairedFZeroDropLo{} a \NThreeBPairedFZeroDropHi),
\NThreeBPairedFOneDrop{} con $F_1$ y \NThreeBPairedFTwoDrop{} con $F_2$, y en
ninguno de los tres casos el intervalo llega a cruzar el cero. El orden de la vecindad reduce la brecha con
independencia de la función de mérito, y lo hace en al menos el
\NThreeBPairedMinBetterPct\,\% de los mundos uno a uno. El efecto de la función de
mérito depende de la regla de revisión en las medianas observadas ---$F_0$
mejora a BR y a LLL y no cambia nada en ASR, con un recorrido de
\NThreeBInteraction{} puntos---, de modo que no se resume mediante un único
efecto marginal.

El factorial compara además las dos familias de reglas. Con el mismo orden, el
mismo criterio de aceptación y los
mismos mundos, la mejor respuesta se queda en \NThreeBBrBestGap\,\% de brecha,
la revisión de Smith atómica en \NThreeBAsrBestGap{} y el aprendizaje
log-lineal en \NThreeBLllBestGap, con factibilidades entre
\NThreeBBrBestFeas{} y \NThreeBLllBestFeas\,\%. La mejor configuración
observada de BR obtuvo una brecha menor que las mejores de ASR y de LLL, aunque
el efecto de la regla depende de la función de mérito: con $F_1$ el orden se
invierte. En su versión continua, las cuatro dinámicas terminan en tasas enteras
muy parecidas y el orden de mérito continuo no anticipa el entero en el
\NFourFiiMismatchPct\,\% de los mundos.

\begin{figure}[htbp]\centering
\pdfpanel{.99\linewidth}{5.6cm}{assets/figures/sp1-final/f9_n3b_factorial.pdf}
{Factorial con brazo de control}
\caption{(A) Brecha mediana de las doce celdas sobre \NThreeBCommon{} mundos de
soporte común; la línea separa el control con $h=2$. (B) Factibilidad
condicionada de las mismas celdas. (C) Diferencia pareada por mundo entre 2BR y
BR con la misma función de mérito, con IC 95\,\% \emph{bootstrap}.}
\label{fig:n3b}
\viuownsource
\end{figure}

El banco industrial reúne \NFourHWorlds{} mundos sobre pasillos cruzados,
pasillos paralelos, muelles con zona de preparación y dos áreas unidas por un
paso estrecho. El oráculo certificó \NFourHCertified; BR entregó salida factible
en \NFourHBrFeasN{} de ellos y 2BR en \NFourHTwoBrFeasN, de modo que el
contraste pareado se calcula sobre los \NFourHPairN{} mundos en que ambas son
factibles. Ahí pasar de $h=1$ a $h=2$ redujo la brecha en
\NFourHPairedDrop{} puntos (IC 95\,\%: \NFourHPairedDropLo{} a
\NFourHPairedDropHi), con mejora en el \NFourHPairedBetterPct\,\% de los
mundos. La comparación entre
los siete métodos de la Figura~\ref{fig:holdout} usa un soporte más estrecho,
los \NFourHCommon{} mundos en que las siete salidas son factibles.

El primer escape conectado siguió siendo bilateral en el
\NFourHHstarTwoPct\,\% de los \NFourHHstarN{} terminales, y el sobrecoste de
DMIS+TX sobre la confirmación global fue de \NFourHDmisOverhead{} bytes/AMR
(IC 95\,\%: \NFourHDmisOverheadLo--\NFourHDmisOverheadHi), del mismo orden que
en el banco propio. Las dos arquitecturas volvieron a coincidir en brecha y en
factibilidad.

La factibilidad, en cambio, disminuye (Figura~\ref{fig:holdout}B): BR resolvió
el \NFourHBrFeas\,\% de los mundos certificados y 2BR el \NFourHTwoBrFeas\,\%,
frente al \SPoneBRFeasRound{} y el \SPoneTwoBRFeasRound\,\% del banco propio, y
CBBA-RB cae al \NFourHCbbaFeas\,\%. El ranking de las medianas se mantiene,
aunque sus valores absolutos empeoran. La dificultad depende de la planta, del
\NFourHEasyGap\,\% de brecha en \NFourHEasyLayout{} al \NFourHHardGap\,\% en
\NFourHHardLayout, y la mayor asociación descriptiva con la brecha se observó para la dispersión
relativa de las distancias de cada AMR a las cargas,
$\overline{(\max_kd_{ik}-\min_kd_{ik})/\overline d_{i}}$: $\rho=\NFourHRhoSpread$
de Spearman calculado sobre los mundos individuales, frente a
\NFourHRhoDegree{} del grado medio y \NFourHRhoLambda{} de $\lambda_2$. El
panel~C de la Figura~\ref{fig:holdout} representa la mediana por geometría, no
los mundos uno a uno. Donde AMR y cargas ocupan bandas
opuestas todas las distancias son parecidas (\NFourHSpreadEasy) y elegir mal
cuesta poco; en pasillos paralelos cada AMR está mucho más cerca de las cargas
de su propio pasillo (\NFourHSpreadHard) y una asignación subóptima se paga
cara. La revisión bilateral generaliza como mejora relativa, pero no compensa
por sí sola las restricciones geométricas de las plantas más estructuradas.

\begin{figure}[htbp]\centering
\pdfpanel{.99\linewidth}{5.8cm}{assets/figures/sp1-final/f10_holdout_industrial.pdf}
{Banco industrial}
\caption{(A) Brecha mediana por planta. (B) Factibilidad condicionada con IC de
Wilson sobre \NFourHWorlds{} mundos. (C) Brecha mediana frente a la dispersión
relativa de $d_{ik}$, una marca por planta; el coeficiente de correlación del
texto se calcula en cambio sobre los mundos individuales.}
\label{fig:holdout}
\viuownsource
\end{figure}

\section{Qué mide cada brecha}

Las distancias al óptimo entero que aparecen en este capítulo no son la misma
magnitud: cada una responde a una restricción distinta
(Figura~\ref{fig:lectura}). La primera es el precio de exigir enteros. El óptimo continuo de N2 queda
\NTwoGapPct\,\% por debajo del entero, y esa diferencia no se recupera con
ningún método: es lo que cuesta que un AMR entre completo o no entre. La
segunda es el precio de limitar el orden de revisión. Con $h=1$ la brecha es
del \SPoneBRGapCommon\,\%, con $h=2$ del \SPoneTwoBRGapCommon{} y con $h=3$ del
\SPoneCThreeGapCommon\,\%, sobre el mismo oráculo y el mismo soporte. Las dos
pérdidas conviven en el mismo eje y tienen causas independientes: una viene de
la atomicidad del AMR y la otra de cuántos pueden revisar a la vez.

Ese orden es, en la práctica, un presupuesto de coordinación. Cada peldaño de
la jerarquía $\mathcal V_1\subseteq\mathcal V_2\subseteq\mathcal V_3$ recorta
la brecha y se paga en comunicación: de \NFourBytesQpgU{} a \NFourBytesQpgP{} y
\NFourBytesCThree{} bytes por AMR. Por eso la comparación entre métodos
distribuidos no tiene un ganador único, sino un compromiso entre calidad y tráfico.

La tercera pérdida es distinta de las anteriores y no se había medido: exigir
que los AMR de una desviación puedan hablarse entre sí. Una mejora puede
existir en el espacio de estrategias y no ser acordable. Sobre
\PoCWorlds{} terminales de la mejor respuesta, reconstruidos desde los mundos
congelados y evaluados con y sin esa exigencia, la restricción afectó al
\PoCWorstPct\,\% de los casos con \PoCWorstRegime{} ---\PoCWorstBlocked{}
mejoras bloqueadas y \PoCWorstDelayed{} desplazadas a un orden mayor--- y a
ninguno en los \PoCZeroRegimes{} regímenes más densos. El precio existe, es
pequeño y decae al densificar el grafo, porque en un grafo de disco los AMR que
ganan intercambiando cargas están cerca, y estar cerca es justo lo que los
hace vecinos.

Por último, el factorial separa dos cosas que suelen confundirse: el protocolo
con que un AMR revisa su decisión y la función de mérito con que la valora.
Cambiar uno no equivale a cambiar la otra, y su efecto no es aditivo.

\begin{figure}[htbp]\centering
\pdfpanel{.99\linewidth}{5.8cm}{assets/figures/sp1-final/f11_lectura_conceptual.pdf}
{Las cuatro magnitudes que separa el capítulo}
\caption{(A) Distancia al óptimo entero del LP relajado y de los tres órdenes de
revisión. El LP no pierde nada: produce una cota inferior optimista porque puede
fraccionar un AMR, de modo que su barra es una brecha de relajación y no una
pérdida del mismo tipo que las de $h$. (B) La misma jerarquía frente a los bytes
por AMR, con los métodos de N3 como referencia; los puntos proceden de campañas
distintas y se leen de forma descriptiva, no como contraste pareado. (C)
Proporción de terminales en que exigir un grupo conexo bloquea o desplaza una
mejora, por densidad del grafo. Las magnitudes comparten eje para facilitar la
lectura y no forman una descomposición aditiva de una misma brecha.}
\label{fig:lectura}
\viuownsource
\end{figure}

\section{Limitaciones}

El banco industrial mide geometría, no operación: las plantas son estáticas y
sin obstáculos dinámicos, y la caída de factibilidad que muestran acota el
dominio en el que las cifras del banco propio pueden leerse. El catálogo de
cargas también es estático, así que la reasignación dinámica y las políticas
aprendidas \citep{dai2024dynamicCoalition,bezerra2025dynamicCoalition}, igual
que el tratamiento de tareas infactibles \citep{shida2025infeasibleTasks},
quedan fuera. Contacto, acoplamiento y reparto de \emph{wrench} quedan para SP2,
y la planificación de trayectorias para SP3. Los mundos de este capítulo son planos y
estáticos. Un reclutamiento verificado aquí no anticipa que el transporte
correspondiente sea ejecutable.

Todas las campañas suponen entrega fiable. Pérdida, retardo, partición transitoria
y fallos bizantinos quedan fuera del modelo, y la partición permanente se empleó
como control negativo. El sobrecoste de \SPoneDmisBytesDeltaCert{} bytes/AMR atribuido a DMIS+TX es nominal y no acota lo que costaría bajo un canal degradado, que se estudia en SP4.

Geo-QPG alcanza mínimos $h$-locales y no certifica optimalidad global. Esa
garantía corresponde al algoritmo exhaustivo de la Proposición~\ref{prop:terminacion},
y las ejecuciones experimentales pueden detenerse por presupuesto de rondas antes
de agotar la vecindad. La búsqueda del primer escape conectado se truncó en $h=3$: en
\NFourHStarTwoN{} y \NFourHStarThreeN{} terminales se identificó ese primer
escape dentro de $h\le3$, y de los \NFourHStarGtThreeN{} restantes solo puede
afirmarse que no apareció ninguna mejora conectada hasta $h=3$. En el diagnóstico N1.E4 el oráculo entero operó con límite de
\NOneMilpTimeLimitSec\,s y dejó \NOneMilpUncertifiedCount{} de
\NOneMilpAudits{} instancias sin certificar; las demás campañas declaran su propio
límite.

En el rango evaluado el coste del refinamiento distribuido crece con $N$, no se
reivindica ninguna mejora asintótica frente al solver central y el análisis de
escala se recoge en el Anexo~\ref{anx:escala}. La alternativa distribuida exacta sí se
midió: DPOP \citep{petcuFaltings2005DPOP} reprodujo el óptimo y la factibilidad
del MILP en \NFourDpopRuns{} de \NFourDpopRuns{} instancias pequeñas
($N\in\{6,8\}$, $K=2$), con una discrepancia máxima de \NFourDpopMaxDeltaPm{} pm. Su mayor tabla de mensajes alcanzó
\NFourDpopMaxTable{} entradas ya en ese tamaño, y ese crecimiento es lo que
impide llevarla a la escala del banco principal. De ahí la elección de la mejora
local acotada. Los bytes de las dos familias continuas excluyen el cierre central
$\mathcal R$ y no son comparables con los de los métodos que deciden sobre
enteros; su convergencia tampoco se afirma, porque falta comprobar la cota de
paso de la discretización. La ventaja del LSAP depende de la heurística voraz con
la que se compara y del orden en que esta recorre las cargas, y $d_{ik}$ es la
distancia euclídea, una cota inferior del recorrido real.

\needspace{5\baselineskip}

% ==========================================================================
% PÁGINA 20/20 · síntesis final
% ==========================================================================
% La sintesis desbordaba unas lineas a una pagina propia. Se toma
% prestado el equivalente a cuatro lineas del margen inferior, que es
% el recurso estandar para una viuda de ese tamano.
\enlargethispage{4\baselineskip}
\section{Síntesis de SP1}


Cada nivel puso a prueba un supuesto con una métrica propia. Contar AMR produjo
en N1 hasta \NOneExtremeFalsePct\,\% de coaliciones que incumplen la capacidad
agregada; el óptimo continuo de N2 quedó \NTwoGapPct\,\% por debajo del entero
sobre \NTwoAtomicityComparable{} pares; y en N3 las variantes GRAPE redujeron la
brecha frente a CBBA-RB a costa de más bytes por AMR.

En N4 la brecha mediana pasó de \SPoneBRGapCommon\,\% con BR a
\SPoneTwoBRGapCommon\,\% con 2BR, y ese salto de orden resistió dos controles:
ninguna de las \NThreeBUnilateralCells{} combinaciones de regla y función de
mérito lo reprodujo con $h=1$ sobre los mismos mundos, y se mantuvo en
\NFourHLayouts{} plantas industriales ajenas al diseño, donde en cambio la
factibilidad cae. DMIS+TX cambió la confirmación global por arbitraje vecinal
por \SPoneDmisBytesDeltaCert{} bytes/AMR, y tras el cierre entero el método
mejor clasificado entre las dinámicas continuas cambió en el
\NFourFiiMismatchPct\,\% de los mundos.

En los bancos evaluados, ampliar la revisión de $h=1$ a $h=2$ reduce de forma
consistente la brecha frente al MILP a cambio de más tráfico, y el efecto se
reproduce en geometrías que no se usaron durante el desarrollo, donde sin
embargo la factibilidad baja y la brecha absoluta sube. El orden $h=2$ se
interpreta por tanto como un compromiso entre calidad y coste de coordinación,
no como una garantía de cercanía uniforme al óptimo.

Las formulaciones continuas alcanzan resultados competitivos tras el cierre
$\mathcal R$ ---el esquema primal-dual queda en \NFourPdClosedGapPct\,\% de
brecha mediana---, pero ese cierre es central y puede alterar el orden de mérito
entre métodos. No constituyen por eso una solución distribuida entera comparable
de extremo a extremo con Geo-QPG, y se mantienen como referencias de relajación.
Lo que este capítulo deja fuera continúa en los siguientes: la formación física
---contacto, acoplamiento y reparto de \emph{wrench}--- pasa a SP2, el movimiento
coordinado a SP3 y la degradación del canal a SP4.

\newpage
\appendix
\setcounter{section}{0}
% \Alph con babel en espanol intercala la enye entre la N y la O. En la
% numeracion de apendices se espera el alfabeto latino, asi que se define
% explicitamente.
\newcommand{\latinalph}[1]{%
  \ifcase#1\or A\or B\or C\or D\or E\or F\or G\or H\or I\or J\or K%
  \or L\or M\or N\or O\or P\or Q\or R\or S\or T\or U\or V\or W%
  \or X\or Y\or Z\fi}
\renewcommand{\thesection}{\latinalph{\value{section}}}
% NOTA: este fichero se incluye desde sp1.tex tras \appendix.
% ==========================================================================
% SP1 · material desplazado fuera del cuerpo del capítulo
%
% Se compila como apendice de sp1.tex mediante \input. Recoge, sin perdida,
% el material que la pasada editorial final retiro del cuerpo. Al integrar el
% capitulo en el TFM cada bloque puede reubicarse en su destino definitivo:
%
%   A. Protocolo Monte Carlo pareado      -> capítulo de metodología
%   B. Pseudocódigo de CBBA-RB y GRAPE    -> anexo de algoritmos
%   C. Taxonomía F-I--F-IV y su tabla     -> anexo de métodos
%   D. Glosario de la familia F-I          -> notación
%   E. Trazabilidad e hipótesis H1/H5     -> conclusiones
%   F. Nota de congelación de semillas    -> anexo de reproducibilidad
%
% Cada bloque conserva el texto exacto que tenía en el cuerpo.
% ==========================================================================

% --------------------------------------------------------------------------
% A. Protocolo Monte Carlo pareado (procedía de SP1, p. 4)

\anxsection{Pseudocódigo de los métodos de N3}{anx:pseudocodigo}
% --------------------------------------------------------------------------
\begin{minipage}[t]{.49\linewidth}
{\scriptsize\textbf{Algoritmo 1. Capacity-CBBA-RB.}\par
\begin{algorithmic}[1]
\Require tabla local $R_i$, barreras $h_i[k]$, versión $v_i$
\State recibir $(k,b,C,v,\tau)$ y descartar versiones antiguas
\State fusionar ganadores por $(\Delta D,-d,-\tau)$
\State reconstruir el prefijo de cada $k$ hasta cubrir $m_k$
\If{$i$ queda fuera del prefijo de $k$}
  \State $h_i[k]\gets\max\{h_i[k],b_{ik}\}$; liberar $k$
\EndIf
\State calcular $r_{ik}=[m_k-\sum_{j\in\widehat C_{ik}}c_j]_+$
\If{$i$ está libre}
  \State elegir $k^\star=\arg\max_k(\min\{c_i,r_{ik}\},-d_{ik},-\tau_i)$
  \State pujar solo si $b_{ik^\star}>h_i[k^\star]$; elevar $v_i$
\Else
  \State conservar el registro vigente
\EndIf
\State difundir cambios; refresco completo cada $N$ activaciones
\Ensure coalición candidata, versiones y bytes emitidos
\end{algorithmic}}
\end{minipage}\hfill
\begin{minipage}[t]{.49\linewidth}
{\scriptsize\textbf{Algoritmo 2. Weighted-/Pair-GRAPE.}\par
\begin{algorithmic}[1]
\Require perfil común $a$, fase $q\in\{0,1\}$, contador local
\State intercambiar $(c_i,d_i[\cdot],v_i)$ y verificar el hash de $a$
\State calcular la mejor desviación que reduzca $(D,J)$
\State proponer $(\Delta D,\Delta J,i,a_i',v_i)$ o propuesta nula
\For{$\ell=1,\ldots,N-2$}
  \State inundar el máximo por orden total; reintentar hasta 4 veces
\EndFor
\State en $\ell=N-1$, aplicar una única propuesta en todos los nodos
\State incrementar época y difundir versión y hash del perfil
\If{no hay propuesta y $q=0$ en Pair-GRAPE}
  \State fijar $q\gets1$; enumerar pares con un vecino
\ElsIf{no hay propuesta}
  \State terminar y verificar coincidencia de perfiles
\Else
  \State volver a la etapa de propuesta
\EndIf
\Ensure perfil terminal común o fallo operacional por desacuerdo
\end{algorithmic}}
\end{minipage}

% --------------------------------------------------------------------------
% C. Taxonomía de familias decisionales de N4 (procedía de SP1, p. 13)

\anxsection{Atlas de reglas y de orden de revisión}{anx:reglas}
% ------------------------------------------------------------------
\begin{figure}[H]\centering
\tikzpanelfit{.98\linewidth}{8.4cm}{figures/n4/n4_order_revision.tex}
{Orden estratégico y regla de revisión}
\caption{Vecindades $\mathcal V_1$, $\mathcal V_2$ y $\mathcal V_3$ frente a las
reglas BR, ASR y LLL.}
\viuderivedsource{la revisión de \citet{sandholm2010population} y el
aprendizaje log-lineal de \citet{mardenShamma2012LogLinear}}
\end{figure}

\begin{figure}[H]\centering
\tikzpanelfit{.98\linewidth}{11.2cm}{figures/n4/n4_algorithm_atlas_fi.tex}
{BR, Geo-ASR y Geo-LLL}
\caption{BR, Geo-ASR y Geo-LLL sobre un mismo estado inicial: regla de selección y
desviación aceptada dentro de $\mathcal V_1(a)$.}
\viuderivedsource{la revisión de \citet{sandholm2010population} y el
aprendizaje log-lineal de \citet{mardenShamma2012LogLinear}}
\end{figure}

\begin{figure}[H]\centering
\tikzpanelfit{.98\linewidth}{11.2cm}{figures/n4/n4_algorithm_atlas_coalitional.tex}
{2BR, C3 y DMIS+TX}
\caption{2BR y C3 sobre el mismo estado: desviaciones de orden dos y tres. DMIS+TX:
grafo de conflictos, conjunto independiente y confirmación transaccional con
versiones.}
\viuderivedsource{\citet{luby1986MIS} y \citet{grayLamport2006Commit}}
\end{figure}

% ------------------------------------------------------------------
% J. Esquema de los comparadores continuos (procedia de SP1, F-II/F-III)

\anxsection{Protocolo experimental y análisis de escala}{anx:escala}
%    Destino: capitulo comun de metodologia.
% ------------------------------------------------------------------
\subsection{Mundos, generadores y contraste pareado}

Los mundos se generan con cinco geometrías normalizadas
(Figura~\ref{fig:geometrias}) y la Tabla~\ref{tab:protocolo} recoge sus
parámetros junto con la regla de inferencia. Cada semilla genera un mundo
independiente. Dentro de un mismo mundo todos los
métodos reciben los mismos AMR, cargas y parámetros, y el \emph{bootstrap}
remuestrea mundos completos por celda. Las ejecuciones fallidas o censuradas se
conservan en los datos, y Holm corrige cada familia de contrastes. Las campañas
siguen el protocolo Monte Carlo pareado definido en el capítulo de metodología.

\begin{figure}[H]\centering
\tikzpanel{.98\linewidth}{9.8cm}{figures/compact/fig03_scenarios.tex}
{Cinco geometrías y retirada de un AMR}
\caption{Cinco generadores espaciales normalizados y el tratamiento con retirada
de un AMR.}
\label{fig:geometrias}
\viuownsource
\end{figure}

\begin{table}[H]
\begin{minipage}[t]{.57\linewidth}
{\footnotesize
\textbf{(a) Generadores espaciales.}\par
\renewcommand{\arraystretch}{1.08}
\begin{tabularx}{\linewidth}{@{}p{1.55cm}p{3.25cm}X@{}}
\toprule
\textbf{Geo.}&\textbf{AMR}&\textbf{Cargas}\\
\midrule
Aleatorio&$U([0,L_x]\!\times\![0,L_y])$&igual\\
Agrupado&uniforme&dos centros + $\mathcal N(0,\Sigma)$\\
Separado&$x\le.35L_x$&$x\ge.65L_x$\\
Anillo&$r\in[.36,.48]L$&$r\in[.04,.22]L$\\
Pasillo&$\sigma_y=.04L_y$&$\sigma_y=.08L_y$\\
\bottomrule
\end{tabularx}
\par\smallskip
Agrupado, Anillo y Pasillo aplican \texttt{clip}; la retirada de un AMR conserva
la geometría Aleatorio.
}
\end{minipage}\hfill
\begin{minipage}[t]{.39\linewidth}
{\footnotesize
\textbf{(b) Regla estadística común.}\par
\renewcommand{\arraystretch}{1.08}
\begin{tabularx}{\linewidth}{@{}>{\raggedright\arraybackslash}p{1.9cm}>{\raggedright\arraybackslash}X@{}}
\toprule
\textbf{Respuesta}&\textbf{Inferencia}\\
\midrule
Binaria&riesgo pareado; \mbox{McNemar}\\
Continua&mediana + IC \emph{bootstrap}; signo\\
Tres o más&Friedman + Kendall $W$\\
Censura&tiempo truncado; \mbox{supervivencia}\\
\bottomrule
\end{tabularx}%
}%
\end{minipage}
\caption{(a) Parámetros de los cinco generadores espaciales. (b) Regla de
inferencia aplicada según el tipo de respuesta.}
\label{tab:protocolo}
\viuownsource
\end{table}

% ------------------------------------------------------------------
% M. Sensibilidad de P(h_c* > 3) (procedia del cuerpo, N4.E9)
% ------------------------------------------------------------------

\anxsection{Sensibilidad del primer escape conectado}{anx:sensibilidad}

Esos \NFourHStarGtThreeN{} casos se concentran en presión alta y capacidad
homogénea. Dentro del bloque confirmatorio, que fija $N/K=\NFourMainBlockNk$, la
proporción sin escape detectado fue del \NFourNoEscRhoHigh\,\% con
$\rho=0{,}85$ y del \NFourNoEscRhoLow\,\% con $\rho=0{,}70$
($n=\NFourNoEscRhoHighN$ por nivel), y del \NFourNoEscCvZero\,\% con
$\mathrm{CV}=0$, mientras que en los tres niveles de dispersión
restantes no pasó del \NFourNoEscCvMid\,\%
($n=\NFourNoEscCvZeroN$ por nivel).

La relación $N/K$ se barre en un bloque de sensibilidad aparte, de
\NFourRatioBlockN{} mundos, así que sus cifras no comparten denominador con las
anteriores: \NFourNoEscNkTwo\,\% con $N/K=2$, \NFourNoEscNkThree\,\% con $N/K=3$ y
\NFourNoEscNkFour\,\% con $N/K=4$ ($n=\NFourNoEscNkTwoN$ por nivel). Dentro de
ese bloque, la tasa sin escape disminuye de forma marcada al pasar de dos a cuatro
AMR por carga; no se compara con CV ni con la presión, que provienen de otro
diseño.

Los dos bloques responden a diseños distintos: el confirmatorio fija
$N/K=\NFourMainBlockNk$ y barre CV y presión, mientras que el de sensibilidad
barre $N/K$. Sus porcentajes no son términos de una misma descomposición.

% ------------------------------------------------------------------
% N. Formulacion de los comparadores continuos (procedia del cuerpo)
% ------------------------------------------------------------------

\anxsection{Formulación de los comparadores continuos}{anx:formulacion}

\begin{minipage}[t]{.49\linewidth}
\textbf{Dinámicas poblacionales}
\citep{sandholm2010population,quijano2017population}.
Cada AMR reparte una unidad de masa entre las $K+1$ opciones, con $x_i$ en el
símplex $\mathcal X_i$. La cobertura es $Q_k(x)=\sum_ic_i^{\rm pay}x_{ik}$, el
déficit relativo $r_k=[1-Q_k(x)/m_k]_+$ y $\overline{d}_{ik}$ la distancia
normalizada. Con ellos:
\begin{equation}
\begin{aligned}
\label{eq:n4-fii-potential}
\widetilde\Phi(x)&=-\frac{\alpha}{2}\sum_kr_k(x)^2
-\frac{\beta}{N}\sum_{i,k}\overline{d}_{ik}x_{ik},\\[-1mm]
F_{ik}&=\partial_{x_{ik}}\widetilde\Phi .
\end{aligned}
\tag{7}
\end{equation}
\begin{equation}
\begin{aligned}
\dot x_{ik}^{\rm Rep}&=x_{ik}(F_{ik}-\overline{F}_i),\\
\dot x_{ik}^{\rm Smith}&=\sum_qx_{iq}[F_{ik}-F_{iq}]_+
\\[-1mm]&\quad-x_{ik}\sum_q[F_{iq}-F_{ik}]_+,\\
\dot x_{ik}^{\rm BNN}&=[F_{ik}-\overline{F}_i]_+
\\[-1mm]&\quad-x_{ik}\sum_q[F_{iq}-\overline{F}_i]_+,\\
\dot x_{ik}^{\rm Logit}&=
\frac{e^{F_{ik}/\tau}}{\sum_qe^{F_{iq}/\tau}}-x_{ik}.
\end{aligned}
\tag{8}
\end{equation}
$\overline{F}_i=\sum_qx_{iq}F_{iq}$ es el pago medio de $i$ y $\overline{d}_{ik}=d_{ik}/\max_{j,q}d_{jq}$.
\end{minipage}\hfill
\begin{minipage}[t]{.49\linewidth}
\textbf{Primal-dual con restricción compartida.}
\begin{equation}
\begin{aligned}
g_k(x)&=1-\frac{Q_k(x)}{m_k}\le0,\\[-1mm]
J_i&=\frac1N\sum_k\overline{d}_{ik}x_{ik}+\frac\epsilon2\|x_i\|^2 .
\end{aligned}
\tag{9}
\end{equation}
\begin{equation}
\begin{aligned}
\min_{x_i\in\mathcal X_i}\quad&J_i(x_i)\\[-1mm]
\mathrm{s.a.}\quad&g(x)\le0,\qquad i=1,\ldots,N .
\end{aligned}
\tag{10}
\end{equation}
\begin{equation}
\begin{aligned}
\dot x_i&=\Pi_{T_{\mathcal X_i}}[-\nabla_{x_i}\mathcal L_i],\\[-1mm]
\dot\lambda_k&=\Pi_{T_{\mathbb R_+}}[g_k(x)] .
\end{aligned}
\tag{11}
\label{eq:n4-primal-dual}
\end{equation}
con $\mathcal L_i=J_i(x_i)+\lambda^{\!\top}g(x)$.
\begin{equation}
\begin{aligned}
\dot\lambda_{ik}&=\Pi_{T_{\mathbb R_+}}
\!\bigl[\widehat g_{ik}\\[-1mm]
&\qquad-\gamma\!\sum_{j\in\mathcal N_i}
(\lambda_{ik}-\lambda_{jk})\bigr].
\end{aligned}
\tag{12}
\label{eq:n4-consensus}
\end{equation}
$\widehat g_{ik}$ es la estimación local que $i$ mantiene de $g_k$.
\begin{equation}
\label{eq:n4-kkt}
\begin{gathered}
0\in\mathbf F(x^\star)+\nabla g(x^\star)^\top\lambda^\star
+N_{\mathcal X}(x^\star),\\
g(x^\star)\le0,\ \lambda^\star\ge0,\ 
{\lambda^\star}^{\!\top}g(x^\star)=0 .
\end{gathered}
\tag{13}
\end{equation}
$\mathbf F(x)=(\nabla_{x_i}J_i(x))_{i}$ es el pseudogradiente y $N_{\mathcal X}$ el
cono normal a $\mathcal X=\prod_i\mathcal X_i$.
\end{minipage}

% ------------------------------------------------------------------
% O. Nota de interfaz SP1 -> SP2 (procedia de la sintesis)
% ------------------------------------------------------------------

\anxsection{Orden conectado frente a mejora}{anx:hstar}
% ------------------------------------------------------------------
\begin{figure}[H]\centering
\pdfpanel{.90\linewidth}{5.2cm}{assets/figures/n4-v4/n4_hstar_gap_link.pdf}
{Orden conectado frente a mejora}
\caption{Mejora de brecha por categoría de $h_c^\star$, sobre los
\NFourHStarLinkN{} mundos en que BR terminó con salida factible y las tres brechas
quedan definidas. Con $h_c^\star=3$ ($n=\NFourHStarLinkThreeN$) y sin escape
detectado ($n=\NFourHStarLinkGtThreeN$) se muestran las observaciones
individuales, no una densidad estimada.}
\viuownsource
\end{figure}

% ------------------------------------------------------------------
% L. Protocolo experimental de SP1 (procedia del cuerpo, p. 3)

% ------------------------------------------------------------------
% destino: documentación de interfaz SP1--SP2
% ------------------------------------------------------------------


\anxsection{Diagnóstico post hoc de conectividad}{anx:conectividad}

Como diagnóstico posterior a la campaña se comprobó si una medida agregada de
densidad basta para ordenar los regímenes. Sobre el área de trabajo $A=L_xL_y$,
la escala de conectividad de los grafos geométricos aleatorios
\citep{penrose2003RandomGeometric} sugiere el índice adimensional
$\Gamma_R=\pi NR^2/(A\log N)$, que aquí no es un umbral, y ese índice no separa
los regímenes. Mundos etiquetados igual ocupan
rangos solapados, porque cuatro de los cinco generadores no reparten los AMR de
forma uniforme y la geometría cambia la densidad efectiva a igualdad de
multiplicador.

La conectividad algebraica tampoco los separa mundo a mundo. Con $L_G$ el
laplaciano del grafo de comunicación, $\lambda_2(L_G)$ ---su segundo autovalor más
pequeño--- crece cuanto más fuertemente unida está la red. Decrece de forma
monótona en mediana al ralear el grafo, pero los rangos de dos regímenes
contiguos se solapan igual que los de $\Gamma_R$. Lo que sí tiene, y
$\Gamma_R$ no, es una propiedad exacta: vale cero si y solo si el grafo está
desconectado. La etiqueta de régimen describe
entonces una consigna de diseño experimental. No mide la conectividad que cada
mundo acaba teniendo.

% ------------------------------------------------------------------
% Q. Preguntas por nivel (procedia del cuerpo; destino: conclusiones)
% ------------------------------------------------------------------

% ------------------------------------------------------------------
% destino: capítulo de conclusiones
% ------------------------------------------------------------------



\end{document}
